% ============================================================
% The Fundamental Theorem of Calculus in Mathlib
% A Dependency-Guided Formal Exposition
% Compile with: pdflatex (twice)
% ============================================================
% ---- verbatim / fancyvrb configuration ----------------------------
% ---- custom abstract: heading + content ---------------------------
% -------------------------------------------------------------------
% ---- subject classification --------------------------------------
% ---- sectioning fonts --------------------------------------------
% ---- column-type tuning parameters -------------------------------
% ---- column types ------------------------------------------------
%% FIXED-MATH: pinned Mathlib commit used throughout the paper.
% ---- title-case preservation -------------------------------------
% ---- spacing inside \maketitle: title / author / abstract --------
% ------------------------------------------------------------------


\documentclass[11pt,reqno]{amsart}
%%%%%%%%%%%%%%%%%%%%%%%%%%%%%%%%%%%%%%%%%%%%%%%%%%%%%%%%%%%%%%%%%%%%%%%%%%%%%%%%%%%%%%%%%%%%%%%%%%%%%%%%%%%%%%%%%%%%%%%%%%%%%%%%%%%%%%%%%%%%%%%%%%%%%%%%%%%%%%%%%%%%%%%%%%%%%%%%%%%%%%%%%%%%%%%%%%%%%%%%%%%%%%%%%%%%%%%%%%%%%%%%%%%%%%%%%%%%%%%%%%%%%%%%%%%%
\usepackage{amsfonts}
\usepackage[T1]{fontenc}
\usepackage[utf8]{inputenc}
\usepackage{titlesec,amsmath,amssymb,amsthm,ragged2e}
\usepackage{array}
\usepackage{graphicx,setspace}
\usepackage[margin=0.75in]{geometry}
\usepackage[colorlinks=true,linkcolor=blue,citecolor=blue,urlcolor=blue]{hyperref}
\usepackage{xcolor}
\usepackage{fancyvrb}

\setcounter{MaxMatrixCols}{10}
%TCIDATA{OutputFilter=LATEX.DLL}
%TCIDATA{Version=5.50.0.2960}
%TCIDATA{<META NAME="SaveForMode" CONTENT="1">}
%TCIDATA{BibliographyScheme=Manual}
%TCIDATA{LastRevised=Tuesday, October 13, 2026 16:49:35}
%TCIDATA{<META NAME="GraphicsSave" CONTENT="32">}

\IfFileExists{tcilatex.tex}{\input{tcilatex}}{}
\fvset{frame=lines,
	framerule=0.7pt,
	rulecolor=\color{blue!60!black},
	framesep=5pt,
	fontsize=\scriptsize,
	xleftmargin=6pt,
	xrightmargin=6pt
}
\RecustomVerbatimEnvironment{verbatim}{Verbatim}{}
\DefineVerbatimEnvironment{verbatimbox}{Verbatim}{frame=single,
	framerule=0.7pt,
	rulecolor=\color{blue!60!black},
	framesep=5pt,
	fontsize=\scriptsize,
	xleftmargin=6pt,
	xrightmargin=6pt
}
\renewcommand{\abstractname}{\fontsize{16}{18}\selectfont Abstract}
\newcommand{\abstractfont}{	\normalfont
	\fontsize{11}{13.2}\selectfont
}
\newcommand{\abstractheadingsep}{\vspace{1em}}
\makeatletter
\renewenvironment{abstract}{	\ifx\maketitle\relax
	\ClassWarning{\@classname}{Abstract should precede
		\protect\maketitle\space in AMS document classes; reported}	\fi
	\global\setbox\abstractbox=\vtop \bgroup
	\abstractfont
	{\centering\abstractname\par}	\abstractheadingsep
	\list{}{\labelwidth\z@
		\leftmargin3pc \rightmargin\leftmargin
		\listparindent\normalparindent \itemindent\z@
		\parsep\z@ \@plus\p@
		\let\fullwidthdisplay\relax
	}	\item\relax
}{	\endlist\egroup
	\ifx\@setabstract\relax \@setabstracta \fi
}
\makeatother
\makeatletter
\renewcommand{\subjclass}[2][]{	\gdef\@subjclass{#2}	\gdef\@subjclassyear{#1}}
\makeatother
\subjclass[2020]{	\textbf{Primary:} 26A42.\,
	\textbf{Secondary:} 28A15; 28B05; 28A25; 68V20; 68V15.}
\keywords{Fundamental Theorem of Calculus, interval integral,
	Bochner integral, Lebesgue integral, Little-o, Mathlib, Lean 4,
	dependency tree, formalisation}
\setcounter{secnumdepth}{3}
\titleformat{\section}
{\normalfont\fontsize{22}{20}\selectfont\bfseries}
{\thesection.}
{0.5em}
{}
\titleformat{\subsection}
{\normalfont\fontsize{18}{16}\selectfont\bfseries}
{\thesubsection}
{0.5em}
{}
\titleformat{\subsubsection}
{\normalfont\fontsize{14}{12}\selectfont\bfseries}
{\thesubsubsection}
{0.5em}
{}
\titlespacing*{\section}{0pt}{3.5ex plus 1ex minus .2ex}{2.3ex plus .2ex}
\titlespacing*{\subsection}{0pt}{3.25ex plus 1ex minus .2ex}{1.5ex plus .2ex}
\titlespacing*{\subsubsection}{0pt}{3.25ex plus 1ex minus .2ex}{1.5ex plus .2ex}
\newcommand{\colLsize}{\scriptsize}
\def\colLbaseline{1.5}
\newcommand{\colLindent}{1.5em}
\newcommand{\colLwidth}{11cm}
\setlength{\arraycolsep}{3pt}
\newcolumntype{T}{>{\centering\arraybackslash}m{1.4cm}}
\newcolumntype{S}{>{\small\centering\arraybackslash}m{2.8cm}}
\newcolumntype{L}{>{\colLsize
		\baselineskip=\colLbaseline\baselineskip
		\setlength{\parindent}{\colLindent}		\RaggedRight\arraybackslash}m{\colLwidth}}
\newcolumntype{K}{>{\footnotesize\centering\arraybackslash}m{2.8cm}}
\newcommand{\lean}[1]{\texttt{#1}}
\newcommand{\tagged}[2]{\hypertarget{#1}{\textbf{#2}}}
\newcommand{\refto}[2]{\hyperlink{#1}{#2}}
\newcommand{\mathlibsha}{\texttt{0c1d5a7e}}
\newcommand{\mathlibsharaw}{0c1d5a7e}
\makeatletter
\let\orig@maketitle\maketitle
\renewcommand{\maketitle}{	\begingroup
	\let\MakeUppercase\@firstofone
	\let\uppercase\@firstofone
	\orig@maketitle
	\endgroup
}
\makeatother
\makeatletter
\newlength{\titleToAuthorSkip}
\newlength{\authorToAbstractSkip}
\setlength{\titleToAuthorSkip}{18pt plus 14pt minus 2pt}
\setlength{\authorToAbstractSkip}{100pt plus 6pt minus 3pt}
\renewcommand{\@maketitle}{	\normalfont\normalsize
	\let\@makefnmark\relax
	\let\@thefnmark\relax
	\@settitle
	\ifx\@empty\@authors\else
	\vspace{\titleToAuthorSkip}	\@setauthors
	\fi
	\ifx\@empty\@date\else
	\vspace{6pt}	\@setdate
	\fi
	\ifx\@empty\@abstract\else
	\vspace{\authorToAbstractSkip}	\fi
	\@setabstract
}
\makeatother
\pagestyle{plain}
\renewcommand{\seriesdefault}{\bfdefault}
\boldmath

\begin{document}
\title{\bfseries{\Large {The Fundamental Theorem of Calculus in Mathlib}}\\
[10pt] {\Large {A Dependency-Guided Formal Exposition}}}
\author{ 
\bfseries{\large \hspace{1.5in} Brahimi Mahdi Tahar\newline
		\text{ \  } \newline
		\text{ \  } \hspace{0.5in}\text{ \  } University of Mohamed Boudiaf of Msila, Algeria\newline
		\text{ \  } \newline
		\href{mailto:mahditahar.brahimi@univ-msila.dz}{mahditahar.brahimi@univ-msila.dz}}%
}

\begin{abstract}
\setstretch{1.5}We present a dependency-guided exposition of the 
Fundamental Theorem of Calculus as formalised in Mathlib. The paper  roots a
depth-5 dependency tree at the one-sided Little-$o$ theorem  \textsf{%
measure\allowbreak\_integral\allowbreak\_sub\allowbreak\_linear\allowbreak
\_is\allowbreak LittleO\allowbreak\_of\allowbreak\_tendsto\allowbreak\_ae}, 
shows that its two-sided lift  \textsf{measure\allowbreak\_integral%
\allowbreak\_sub\allowbreak\_integral\allowbreak
\_sub\allowbreak\_linear\allowbreak\_is\allowbreak
LittleO\allowbreak\_of\allowbreak \_tendsto\allowbreak\_ae} subsumes the
classical derivative form and the  two one-sided variants, and inspects the
resulting axiom closure. The  closure is confined to \textsf{propext}, 
\textsf{funext} and  \textsf{Quot\allowbreak.sound}; \textsf{%
Classical\allowbreak.choice} is  not reached, so the theorem is choice-free.
A Lean~4 proof skeleton is  type-checked against Mathlib commit \texttt{%
0c1d5a7e}{}, and its  \textsf{\#print\ axioms} output is reported verbatim.
The paper  complements the corpus-scale axiom-profile measurement of
Close~(2026)  by exhibiting a concrete, mathematically significant theorem
that sits  inside the choice-free corner of Mathlib.
\end{abstract}

\maketitle

\thispagestyle{empty}

\newpage \pagestyle{plain} \pagenumbering{arabic} \setstretch{1.5}

\section{Introduction}

\label{sec:introduction}

\noindent Mathlib \cite{mathlib4} formalises the Fundamental Theorem of
Calculus not as a single declaration but as a small hierarchy. Its most
general member is the Little-$o$ statement \textsf{measure\allowbreak%
\_integral\allowbreak\_sub\allowbreak\_integral\allowbreak
\_sub\allowbreak\_linear\allowbreak\_is\allowbreak
LittleO\allowbreak\_of\allowbreak \_tendsto\allowbreak\_ae}, from which the
classical derivative form \textsf{integral\allowbreak\_eq\allowbreak\_sub%
\allowbreak\_of\allowbreak\_has\allowbreak Deriv\allowbreak At} and the
one-sided variants \textsf{\_left} and \textsf{\_right} are derived by
direct reduction. The hierarchy is stable across Mathlib releases and is
cited in downstream work on change of variables, Barrow's rule, and the
Bochner integral.

\smallskip\noindent\textbf{What this paper does.} We make three
contributions.

\begin{enumerate}
\item \emph{A dependency tree.} We expand the one-sided Little-$o$  theorem
to depth~5 using the MathlibGraph edge dataset  \cite{mathlibgraph2026,
mathlibgraph, network-structure-mathlib}. The  printed excerpt is chosen so
that all three kernel axioms reached by  the closure --- \textsf{funext}, 
\textsf{Quot\allowbreak.sound} and  \textsf{propext} --- appear inside it
(Section~\ref{sec:tree}).

\item \emph{Statement tables and a proof strategy.} We collect the  fourteen
definitions, thirteen lemmas, thirteen theorems, two  propositions and two
corollaries that make up the hierarchy  (Sections~\ref{sec:definitions}--\ref%
{sec:statements}), verify each  signature against the Mathlib4 documentation
browser  \cite{mathlib4-docs} and the source file \cite{mathlib4-fundthm},
and  give the five-step reduction from the one-sided to the two-sided 
statement (Section~\ref{sec:proofs}).

\item \emph{A verified Lean skeleton.} We give a Lean~4 proof skeleton  for
the two-sided statement that type-checks against Mathlib commit  \texttt{%
0c1d5a7e}{} and, together with its \textsf{\#print\ axioms} output,  shows
that the closure remains confined to \textsf{propext} and  \textsf{%
Quot\allowbreak.sound} (Section~\ref{sec:lean-skeleton}).
\end{enumerate}

\smallskip\noindent\textbf{Choice-freeness.} The absence of \textsf{%
Classical\allowbreak.choice} from the axiom closure of the FTC-2 hierarchy
is not an accident of the particular proof term we examined. A corpus-scale
study of the axiom profile of 707{,}053 Mathlib declarations \cite%
{close2026axiom} shows that only $23$ distinct axiom profiles occur across
the entire library and that \textsf{Classical\allowbreak.choice} reaches
only about half the corpus ($50.84\%$), through a narrow frontier of $206$
declarations that use the axiom directly. Theorems whose closure lies
entirely inside $\{\mathsf{propext},\mathsf{Quot\allowbreak.sound}\}$ --- as
is the case for the FTC-2 statements examined here --- are common in Mathlib
and arise whenever the underlying construction does not require a
nonconstructive selection principle (Section~\ref{sec:axioms}).

\smallskip\noindent\textbf{Organisation.} Section~\ref{sec:axioms} recalls
the three kernel axioms and reports the \textsf{\#print\ axioms} output for
the root declaration. Section~\ref{sec:definitions} collects the
definitions, Section~\ref{sec:statements} the lemmas, theorems, propositions
and corollaries, and Section~\ref{sec:proofs} the five-step reduction.
Section~\ref{sec:scripts} reports the depth-5 dependency tree, the Lean
skeleton, and the generator script. Section~\ref{sec:tables} gives the
timeline, classifications and applications. Section~\ref{sec:conclusion}
summarises.

\section{Axioms}

\label{sec:axioms}

\noindent Lean~4's kernel contains exactly three primitive logical axioms 
\cite{mathlib4, mathlib4-docs}: \textsf{propext}, \textsf{%
Quot\allowbreak.sound}, and \textsf{Classical\allowbreak.choice}. Mathlib
conventionally also reports \textsf{funext} alongside them, although it is
derivable from \textsf{Quot\allowbreak.sound}. \noindent A \textsf{\#print\
axioms} check on the root declaration of the depth-5 tree (the one-sided
theorem \textsf{interval\allowbreak Integral.measure\allowbreak
\_integral\allowbreak \_sub\allowbreak \_linear\allowbreak \_is\allowbreak
LittleO\allowbreak \_of\allowbreak \_tendsto\allowbreak \_ae}) gives the
following output (verbatim, from a local Lean~4 session at Mathlib commit 
\texttt{0c1d5a7e}{}):

\text{ \ } \setstretch{1}\text{ \ } 
\begin{verbatimbox}
#print axioms
  intervalIntegral.measure_integral_sub_linear_isLittleO_of_tendsto_ae
-- 'intervalIntegral.measure_integral_sub_linear_isLittleO_of_tendsto_ae'
--   depends on axioms: [propext, Quot.sound]
--
-- (funext is derivable from Quot.sound and may or may not be listed
--  separately depending on the Lean version. Classical.choice is not
--  reported: the proof is constructive in the choice parameter. The
--  two-sided lift has the same axiom closure as the one-sided statement.)
\end{verbatimbox}
\text{ \ } \setstretch{1.5}\text{ \ } \noindent The same conclusion is
reached independently by the depth-5 tree traversal reported in Section~\ref%
{sec:tree}: \textsf{propext}, \textsf{funext} and \textsf{%
Quot\allowbreak.sound} are the only kernel axioms reached, and \textsf{%
Classical\allowbreak.choice} is not. 
%% FIXED-MATH: choice-freeness of FTC-2, with a corpus-scale reference.
\noindent\textbf{Remark (choice-freeness).} The absence of \textsf{%
Classical\allowbreak.choice} from the axiom closure of \hyperlink{thm:root}{%
Thm\,1} and \hyperlink{thm:linear}{Thm\,2} is not an artefact of the
particular proof term examined here: it is a property of the whole FTC-2
hierarchy as formalised in Mathlib. A corpus-scale study of the axiom
profile of 707{,}053 Mathlib declarations \cite{close2026axiom} shows that
only $23$ distinct axiom profiles occur across the entire library and that 
\textsf{Classical\allowbreak.choice} reaches only about half the corpus ($%
50.84\%$), through a narrow frontier of $206$ declarations that use the
axiom directly. Theorems whose closure lies entirely inside $\{\mathsf{%
propext},\mathsf{Quot\allowbreak.sound}\}$ --- as is the case for the FTC-2
statements examined here --- are common in Mathlib and arise whenever the
underlying construction does not require a nonconstructive selection
principle. In the present case, none of the five reduction steps of Section~%
\ref{sec:proofs} invokes a choice principle: the lemmas \hyperlink{lem:addadd}%
{L1}, \hyperlink{lem:negleft}{L2}, \hyperlink{lem:negright}{L2$^{\prime }$}, 
\hyperlink{lem:compTendsto}{L11} and the interval-integral lemmas \hyperlink{thm:addadj}%
{Thm\,7}--\hyperlink{thm:toIxx}{Thm\,13} are all themselves choice-free, and
the composite inherits their axiom profile.\text{ \ }\newline

\noindent $%
\begin{array}{|T|S|L|K|}
\hline
\hypertarget{ax:propext}{\textbf{Ax\,1}} &
propositional extensionality &
\textsf{propext}$:\forall (P:\textsf{Prop})\;(Q:\textsf{Prop}),\;(P\leftrightarrow Q)\Longrightarrow P=Q$ &
\textsf{ker\allowbreak nel ax\allowbreak iom} \\ \hline
\hypertarget{ax:funext}{\textbf{Ax\,2}} &
function extensionality &
\textsf{funext}$:\forall (\alpha:\textsf{Type})\;(\beta:\textsf{Type})\;(f:\alpha\to\beta)\;(g:\alpha\to\beta),\;(\forall (x:\alpha),\;f(x)=g(x))\Longrightarrow f=g$ &
\textsf{kernel axiom} \\ 
\end{array}$ \noindent $%
\begin{array}{|T|S|L|K|}
\hline
\hypertarget{ax:quot}{\textbf{Ax\,3}} &
quo\allowbreak tient sound\allowbreak ness &
\textsf{Quot\allowbreak.sound}$:\forall (\alpha:\textsf{Type})\;(r:\alpha\to\alpha\to\textsf{Prop})\;(a:\alpha)\;(b:\alpha),\;r(a,b)\Longrightarrow\textsf{Quot\allowbreak.mk}(r,a)=\textsf{Quot\allowbreak.mk}(r,b)$ &
\textsf{ker\allowbreak nel ax\allowbreak iom} \\ \hline
\hypertarget{ax:choice}{\textbf{Ax\,4}} &
ax\allowbreak iom of choice &
\textsf{Classical\allowbreak.choice}$:\forall (\alpha:\textsf{Type}),\;\textsf{Non\allowbreak empty}(\alpha)\Longrightarrow\alpha$ &
\textsf{ker\allowbreak nel ax\allowbreak iom} \\ \hline
\end{array}$ \noindent \text{ \ }\newline

\noindent A full depth-5 expansion of the one-sided theorem (see Section~\ref%
{sec:tree}) reaches all three of \textsf{propext}, \textsf{funext} and 
\textsf{Quot\allowbreak.sound}: \textsf{funext} and \textsf{%
Quot\allowbreak.sound} appear inside the printed excerpt (under \textsf{%
Asymptotics\allowbreak.IsLittleO.neg\allowbreak \_left}), while \textsf{%
propext} is reached on the \textsf{Asymptotics\allowbreak.IsLittleO.sub}
branch immediately outside it. \textsf{Classical\allowbreak.choice} does not
appear anywhere in the depth-5 closure of the tree examined here. The kernel
axioms and their status are documented in \cite{mathlib4-docs}.

\section{Definitions}

\label{sec:definitions}

\noindent The proof of the general FTC-2 rests on the following definitions,
all drawn from the depth-1 frontier of the dependency tree generated from
the MathlibGraph dataset \cite{mathlibgraph2026, mathlibgraph}.

\text{\ } \newline
\text{\ }

\noindent$%
\begin{array}{|T|S|L|K|}
\hline
\hypertarget{def:II}{\textbf{Def\,1}} &
\textsf{Interval\allowbreak Integrable} &
\textsf{Interval\allowbreak Integrable}$:\forall (E:\textsf{Type})\;[Normed\allowbreak AddCommGroup\;E]\;[Normed\allowbreak Space\;\mathbb{R}\;E]\;(f:\mathbb{R}\to E)\;(a:\mathbb{R})\;(b:\mathbb{R})\;(\mu:\textsf{Measure}\;\mathbb{R}),\;\textsf{Prop}$ &
\textsf{Measure\allowbreak Theory} \\ \hline
\hypertarget{def:IIsymm}{\textbf{Def\,2}} &
sym\allowbreak metry &
\textsf{Interval\allowbreak Integrable.symm}$:\forall (E:\textsf{Type})\;[Normed\allowbreak AddCommGroup\;E]\;[Normed\allowbreak Space\;\mathbb{R}\;E]\;(f:\mathbb{R}\to E)\;(a:\mathbb{R})\;(b:\mathbb{R})\;(\mu:\textsf{Measure}\;\mathbb{R}),\;\textsf{Interval\allowbreak Integrable}(E,f,a,b,\mu)\Longrightarrow\textsf{Interval\allowbreak Integrable}(E,f,b,a,\mu)$ &
\textsf{Measure\allowbreak Theory} \\ 
\end{array}$ \noindent $%
\begin{array}{|T|S|L|K|}
\hline
\hypertarget{def:IItrans}{\textbf{Def\,3}} &
transi\allowbreak tivity &
\textsf{Interval\allowbreak Integrable.trans}$:\forall (E:\textsf{Type})\;[Normed\allowbreak AddCommGroup\;E]\;[Normed\allowbreak Space\;\mathbb{R}\;E]\;(f:\mathbb{R}\to E)\;(a:\mathbb{R})\;(b:\mathbb{R})\;(c:\mathbb{R})\;(\mu:\textsf{Measure}\;\mathbb{R}),\;\textsf{Interval\allowbreak Integrable}(E,f,a,b,\mu)\Longrightarrow\textsf{Interval\allowbreak Integrable}(E,f,b,c,\mu)\Longrightarrow\textsf{Interval\allowbreak Integrable}(E,f,a,c,\mu)$ &
\textsf{Measure\allowbreak Theory} \\ \hline
\hypertarget{def:FTCF}{\textbf{Def\,4}} &
FTC filter type\allowbreak class &
\textsf{FTC\allowbreak Filter}$:\forall (\ell:\textsf{Filter}\;\mathbb{R})\;(a:\mathbb{R}),\;\textsf{Prop}$ &
\textsf{Interval\allowbreak Integral} \\ \hline
\end{array}$

\text{\ } \newline
\text{\ } \noindent\textit{Note on \hyperlink{def:FTCF}{Def 4}.} \textsf{%
FTC\allowbreak Filter} is a two-argument predicate \textsf{FTC\allowbreak
Filter} $:\mathsf{Filter}\;\mathbb{R}\to\mathbb{R}\to\mathsf{Prop}$, written 
\textsf{[FTC\allowbreak Filter}$\;\ell\;a$\textsf{]} in Lean~4 syntax. The
brackets denote that it is used as a typeclass instance. This is the \emph{%
canonical} notation used throughout the paper: the same spelling \textsf{%
[FTC\allowbreak Filter}$\;\ell\;a$\textsf{]} appears in the statement tables
of Section~\ref{sec:statements} and in the verbatim Lean skeleton of Section~%
\ref{sec:lean-skeleton}. \text{\ } \newline
\text{\ }

\noindent$%
\begin{array}{|T|S|L|K|}
\hline
\hypertarget{def:TIC}{\textbf{Def\,5}} &
\textsf{Tendsto\allowbreak Ixx\allowbreak Class} &
\textsf{Tendsto\allowbreak Ixx\allowbreak Class}$:\forall (\alpha:\textsf{Type})\;(\beta:\textsf{Type})\;(F:\textsf{Filter}\;\alpha)\;(\ell:\textsf{Filter}\;\beta)\;(\textsf{Ixx}:\beta\to\beta\to\textsf{Set}\;\alpha),\;\textsf{Prop}$ &
\textsf{Filter} \\ \hline
\hypertarget{def:IMG}{\textbf{Def\,6}} &
measur\allowbreak ably gener\allowbreak ated &
\textsf{Is\allowbreak Measurably\allowbreak Generated}$:\forall (\alpha:\textsf{Type})\;[\textsf{Measurable\allowbreak Space}\;\alpha]\;(\ell:\textsf{Filter}\;\alpha),\;\textsf{Prop}$ &
\textsf{Filter} \\ 
\end{array}$

\noindent$%
\begin{array}{|T|S|L|K|}
\hline
\hypertarget{def:IsLO}{\textbf{Def\,7}} &
Lit\allowbreak tle-$o$ rela\allowbreak tion &
\textsf{Is\allowbreak LittleO}$:\forall (\alpha:\textsf{Type})\;(E:\textsf{Type})\;(F:\textsf{Type})\;[Normed\allowbreak AddCommGroup\;E]\;[Normed\allowbreak AddCommGroup\;F]\;(\ell:\textsf{Filter}\;\alpha)\;(f:\alpha\to E)\;(g:\alpha\to F),\;\textsf{Prop}$ &
\textsf{Asymp\allowbreak totics} \\ \hline
\hypertarget{def:Tendsto}{\textbf{Def\,8}} &
filter conver\allowbreak gence &
\textsf{Tendsto}$:\forall (\alpha:\textsf{Type})\;(\beta:\textsf{Type})\;(f:\alpha\to\beta)\;(\ell_{1}:\textsf{Filter}\;\alpha)\;(\ell_{2}:\textsf{Filter}\;\beta),\;\textsf{Prop}$ &
\textsf{Filter} \\ 
\end{array}$

\noindent$%
\begin{array}{|T|S|L|K|}
\hline
\hypertarget{def:NACG}{\textbf{Def\,9}} &
nor\allowbreak med addi\allowbreak tive comm.\ group &
\textsf{Normed\allowbreak AddCommGroup}$:\forall (E:\textsf{Type}),\;\textsf{Prop}$ &
\textsf{Anal\allowbreak ysis.\allowbreak Normed} \\ \hline
\hypertarget{def:NS}{\textbf{Def\,10}} &
nor\allowbreak med space over $\mathbb{R}$ &
\textsf{Normed\allowbreak Space}$:\forall (\mathbb{K}:\textsf{Type})\;(V:\textsf{Type})\;[Normed\allowbreak Field\;\mathbb{K}]\;[Seminormed\allowbreak AddCommGroup\;V],\;\textsf{Prop}$ &
\textsf{Anal\allowbreak ysis.\allowbreak Normed} \\ 
\end{array}$

\noindent$%
\begin{array}{|T|S|L|K|}
\hline
\hypertarget{def:IPC}{\textbf{Def\,11}} &
forget\allowbreak ful coercion &
\textsf{Inner\allowbreak Product\allowbreak Space.to\allowbreak Normed\allowbreak Space}$:\forall (\mathbb{K}:\textsf{Type})\;(E:\textsf{Type})\;[Normed\allowbreak Field\;\mathbb{K}]\;[Seminormed\allowbreak AddCommGroup\;E],\;\textsf{Inner\allowbreak Product\allowbreak Space}(\mathbb{K},E)\Longrightarrow\textsf{Normed\allowbreak Space}(\mathbb{K},E)$ &
\textsf{Inner\allowbreak Product\allowbreak Space} \\ \hline
\hypertarget{def:LFM}{\textbf{Def\,12}} &
loc\allowbreak ally fin\allowbreak ite mea\allowbreak sure &
\textsf{Is\allowbreak Locally\allowbreak Finite\allowbreak Measure}$:\forall (\alpha:\textsf{Type})\;[\textsf{Measurable\allowbreak Space}\;\alpha]\;[\textsf{Topological\allowbreak Space}\;\alpha]\;(\mu:\textsf{Measure}\;\alpha),\;\textsf{Prop}$ &
\textsf{Measure\allowbreak Theory} \\ 
\end{array}$

\noindent$%
\begin{array}{|T|S|L|K|}
\hline
\hypertarget{def:vol}{\textbf{Def\,13}} &
canon\allowbreak ical mea\allowbreak sure &
\textsf{volume}$:\forall (\alpha:\textsf{Type})\;[\textsf{Measurable\allowbreak Space}\;\alpha],\;\textsf{Measure}\;\alpha$ &
\textsf{Measure\allowbreak Theory} \\ \hline
\hypertarget{def:uIcc}{\textbf{Def\,14}} &
unord\allowbreak ered inter\allowbreak val &
\textsf{uIcc}$:\forall (\alpha:\textsf{Type})\;[\textsf{Linear\allowbreak Order}\;\alpha]\;(a:\alpha)\;(b:\alpha),\;\textsf{Set}\;\alpha$ &
\textsf{Data.\allowbreak Set} \\ \hline
\end{array}$ \noindent \setstretch{1}

\section{Statements}

\label{sec:statements} \setstretch{1.3} \noindent The declarations involved
in the general FTC-2 fall into four logical categories. Every signature
below is taken from the Mathlib4 documentation browser \cite{mathlib4-docs}
and cross-checked against the source file \cite{mathlib4-fundthm} at commit 
\texttt{0c1d5a7e}{}. \setstretch{1}

\subsection{Lemmas}

\label{sec:lemmas} \setstretch{1.4} \noindent Small technical lemmas invoked
inside the proof.

\text{ \ }\newline
\text{ \ } \noindent$%
\begin{array}{|T|S|L|K|}
\hline
\hypertarget{lem:addadd}{\textbf{Lem\,1}} &
sum of two Lit\allowbreak tle-$o$ &
\textsf{Is\allowbreak LittleO.add\allowbreak\_add}$:\forall (\alpha:\textsf{Type})\;(E:\textsf{Type})\;(F:\textsf{Type})\;[Normed\allowbreak AddCommGroup\;E]\;[Normed\allowbreak AddCommGroup\;F]\;(\ell:\textsf{Filter}\;\alpha)\;(f_{1}:\alpha\to E)\;(f_{2}:\alpha\to E)\;(g_{1}:\alpha\to F)\;(g_{2}:\alpha\to F),\;\textsf{Is\allowbreak LittleO}(\alpha,E,F,\ell,f_{1},g_{1})\Longrightarrow\textsf{Is\allowbreak LittleO}(\alpha,E,F,\ell,f_{2},g_{2})\Longrightarrow\textsf{Is\allowbreak LittleO}(\alpha,E,F,\ell,f_{1}+f_{2},g_{1}+g_{2})$ &
\textsf{Asymp\allowbreak totics} \\ \hline
\hypertarget{lem:negleft}{\textbf{Lem\,2}} &
nega\allowbreak tion on the left &
\textsf{Is\allowbreak LittleO.neg\allowbreak\_left}$:\forall (\alpha:\textsf{Type})\;(E:\textsf{Type})\;(F:\textsf{Type})\;[Normed\allowbreak AddCommGroup\;E]\;[Normed\allowbreak AddCommGroup\;F]\;(\ell:\textsf{Filter}\;\alpha)\;(f:\alpha\to E)\;(g:\alpha\to F),\;\textsf{Is\allowbreak LittleO}(\alpha,E,F,\ell,f,g)\Longrightarrow\textsf{Is\allowbreak LittleO}(\alpha,E,F,\ell,-f,g)$ &
\textsf{Asymp\allowbreak totics} \\ 
\end{array}$

\noindent$%
\begin{array}{|T|S|L|K|}
\hline
\hypertarget{lem:negright}{\textbf{Lem\,2$'$}} &
nega\allowbreak tion on the right &
\textsf{Is\allowbreak LittleO.neg\allowbreak\_right}$:\forall (\alpha:\textsf{Type})\;(E:\textsf{Type})\;(F:\textsf{Type})\;[Normed\allowbreak AddCommGroup\;E]\;[Normed\allowbreak AddCommGroup\;F]\;(\ell:\textsf{Filter}\;\alpha)\;(f:\alpha\to E)\;(g:\alpha\to F),\;\textsf{Is\allowbreak LittleO}(\alpha,E,F,\ell,f,g)\Longrightarrow\textsf{Is\allowbreak LittleO}(\alpha,E,F,\ell,f,-g)$ &
\textsf{Asymp\allowbreak totics} \\ \hline
\hypertarget{lem:congr'}{\textbf{Lem\,3}} &
congru\allowbreak ence both argu\allowbreak ments &
\textsf{Is\allowbreak LittleO.congr'}$:\forall (\alpha:\textsf{Type})\;(E:\textsf{Type})\;(F:\textsf{Type})\;[Normed\allowbreak AddCommGroup\;E]\;[Normed\allowbreak AddCommGroup\;F]\;(\ell:\textsf{Filter}\;\alpha)\;(f_{1}:\alpha\to E)\;(f_{2}:\alpha\to E)\;(g_{1}:\alpha\to F)\;(g_{2}:\alpha\to F),\;f_{1}=^{f}[\ell]\;f_{2}\Longrightarrow g_{1}=^{f}[\ell]\;g_{2}\Longrightarrow(\textsf{Is\allowbreak LittleO}(\alpha,E,F,\ell,f_{1},g_{1})\leftrightarrow\textsf{Is\allowbreak LittleO}(\alpha,E,F,\ell,f_{2},g_{2}))$ &
\textsf{Asymp\allowbreak totics} \\ 
\end{array}$

\noindent$%
\begin{array}{|T|S|L|K|}
\hline
\hypertarget{lem:congrL}{\textbf{Lem\,4}} &
congru\allowbreak ence left argu\allowbreak ment &
\textsf{Is\allowbreak LittleO.congr\allowbreak\_left}$:\forall (\alpha:\textsf{Type})\;(E:\textsf{Type})\;(F:\textsf{Type})\;[Normed\allowbreak AddCommGroup\;E]\;[Normed\allowbreak AddCommGroup\;F]\;(\ell:\textsf{Filter}\;\alpha)\;(f_{1}:\alpha\to E)\;(f_{2}:\alpha\to E)\;(g:\alpha\to F),\;\textsf{Is\allowbreak LittleO}(\alpha,E,F,\ell,f_{1},g)\Longrightarrow f_{1}=^{f}[\ell]\;f_{2}\Longrightarrow\textsf{Is\allowbreak LittleO}(\alpha,E,F,\ell,f_{2},g)$ &
\textsf{Asymp\allowbreak totics} \\ \hline
\hypertarget{lem:transle}{\textbf{Lem\,5}} &
trans\allowbreak fer across a bound &
\textsf{Is\allowbreak LittleO.trans\allowbreak\_le}$:\forall (\alpha:\textsf{Type})\;(E:\textsf{Type})\;(F:\textsf{Type})\;[Normed\allowbreak AddCommGroup\;E]\;[Normed\allowbreak AddCommGroup\;F]\;(\ell:\textsf{Filter}\;\alpha)\;(f:\alpha\to E)\;(g:\alpha\to F)\;(h:\alpha\to F),\;\textsf{Is\allowbreak LittleO}(\alpha,E,F,\ell,f,g)\Longrightarrow(\forall^{f}x\;\textsf{in}\;\ell,\;\lVert g(x)\rVert\le\lVert h(x)\rVert)\Longrightarrow\textsf{Is\allowbreak LittleO}(\alpha,E,F,\ell,f,h)$ &
\textsf{Asymp\allowbreak totics} \\ 
\end{array}$

\noindent$%
\begin{array}{|T|S|L|K|}
\hline
\hypertarget{lem:isLOiff}{\textbf{Lem\,6}} &
Lit\allowbreak tle-$o$ via $\varepsilon$ &
\textsf{is\allowbreak LittleO\allowbreak\_iff}$:\forall (\alpha:\textsf{Type})\;(E:\textsf{Type})\;(F:\textsf{Type})\;[Normed\allowbreak AddCommGroup\;E]\;[Normed\allowbreak AddCommGroup\;F]\;(\ell:\textsf{Filter}\;\alpha)\;(f:\alpha\to E)\;(g:\alpha\to F),\;\textsf{Is\allowbreak LittleO}(\alpha,E,F,\ell,f,g)\leftrightarrow\forall (c:\mathbb{R}),\;c>0\Longrightarrow\forall^{f}x\;\textsf{in}\;\ell,\;\lVert f(x)\rVert\le c\cdot\lVert g(x)\rVert$ &
\textsf{Asymp\allowbreak totics} \\ \hline
\hypertarget{lem:mpmem}{\textbf{Lem\,7}} &
member\allowbreak ship under super\allowbreak set &
\textsf{Filter.mp\allowbreak\_mem}$:\forall (\alpha:\textsf{Type})\;(\ell:\textsf{Filter}\;\alpha)\;(s:\textsf{Set}\;\alpha)\;(t:\textsf{Set}\;\alpha),\;\ell(s)\Longrightarrow(\forall (x:\alpha),\;s(x)\Longrightarrow t(x))\Longrightarrow\ell(t)$ &
\textsf{Filter} \\ 
\end{array}$

\noindent$%
\begin{array}{|T|S|L|K|}
\hline
\hypertarget{lem:univmem}{\textbf{Lem\,8}} &
\textsf{univ} in every filter &
\textsf{Filter.univ\allowbreak\_mem'}$:\forall (\alpha:\textsf{Type})\;(\ell:\textsf{Filter}\;\alpha),\;\ell(\textsf{univ})$ &
\textsf{Filter} \\ \hline
\hypertarget{lem:constpure}{\textbf{Lem\,9}} &
con\allowbreak stant tends to pure &
\textsf{Filter.tendsto\allowbreak\_const\allowbreak\_pure}$:\forall (\alpha:\textsf{Type})\;(a:\alpha)\;(\ell:\textsf{Filter}\;\alpha),\;\textsf{Tendsto}(\alpha,\alpha,(\lambda\_:\alpha,\;a),\ell,\textsf{pure}(a))$ &
\textsf{Filter} \\ 
\end{array}$

\noindent$%
\begin{array}{|T|S|L|K|}
\hline
\hypertarget{lem:monoright}{\textbf{Lem\,10}} &
monoton\allowbreak icity on the right &
\textsf{Filter.Tendsto.mono\allowbreak\_right}$:\forall (\alpha:\textsf{Type})\;(\beta:\textsf{Type})\;(f:\alpha\to\beta)\;(\ell_{1}:\textsf{Filter}\;\alpha)\;(\ell_{2}:\textsf{Filter}\;\beta)\;(\ell_{3}:\textsf{Filter}\;\beta),\;\textsf{Tendsto}(\alpha,\beta,f,\ell_{1},\ell_{2})\Longrightarrow\ell_{2}\le\ell_{3}\Longrightarrow\textsf{Tendsto}(\alpha,\beta,f,\ell_{1},\ell_{3})$ &
\textsf{Filter} \\ \hline
\hypertarget{lem:compTendsto}{\textbf{Lem\,11}} &
compo\allowbreak sition with conver\allowbreak gence &
\textsf{Is\allowbreak LittleO.comp\allowbreak\_tendsto}$:\forall (\alpha:\textsf{Type})\;(\beta:\textsf{Type})\;(E:\textsf{Type})\;(F:\textsf{Type})\;[Normed\allowbreak AddCommGroup\;E]\;[Normed\allowbreak AddCommGroup\;F]\;(f:\beta\to E)\;(g:\beta\to F)\;(u:\alpha\to\beta)\;(\ell:\textsf{Filter}\;\alpha)\;(\ell':\textsf{Filter}\;\beta),\;\textsf{Is\allowbreak LittleO}(\beta,E,F,\ell',f,g)\Longrightarrow\textsf{Tendsto}(\alpha,\beta,u,\ell,\ell')\Longrightarrow\textsf{Is\allowbreak LittleO}(\alpha,E,F,\ell,f\circ u,g\circ u)$ &
\textsf{Asymp\allowbreak totics} \\ 
\end{array}$

\noindent$%
\begin{array}{|T|S|L|K|}
\hline
\hypertarget{lem:forallCongr}{\textbf{Lem\,12}} &
$\forall$ congru\allowbreak ence &
\textsf{forall\allowbreak\_congr}$:\forall (\alpha:\textsf{Type})\;(p:\alpha\to\textsf{Prop})\;(q:\alpha\to\textsf{Prop}),\;(\forall (x:\alpha),\;p(x)\leftrightarrow q(x))\Longrightarrow(\forall (x:\alpha),\;p(x))\leftrightarrow(\forall (x:\alpha),\;q(x))$ &
\textsf{quan\allowbreak tifier} \\ \hline
\hypertarget{lem:impliesCongr}{\textbf{Lem\,13}} &
$\to$ congru\allowbreak ence &
\textsf{implies\allowbreak\_congr}$:\forall (p:\textsf{Prop})\;(p':\textsf{Prop})\;(q:\textsf{Prop})\;(q':\textsf{Prop}),\;(p\leftrightarrow p')\Longrightarrow(q\leftrightarrow q')\Longrightarrow((p\to q)\leftrightarrow(p'\to q'))$ &
\textsf{quan\allowbreak tifier} \\ \hline
\end{array}$

\subsection{Theorems}

\label{sec:theorems} The substantive results in the dependency tree. The
root theorem and its one-sided reduction are taken from \cite%
{mathlib4-fundthm} at commit \texttt{0c1d5a7e}{}.

\text{ \ }\newline
\noindent$\setstretch{1.3} 
\begin{array}{|T|S|L|K|}
\hline
\hypertarget{thm:root}{\textbf{Thm\,1}} &
two-\allowbreak sided Lit\allowbreak tle-$o$ FTC-2 &
\textsf{measure\allowbreak\_integral\allowbreak\_sub\allowbreak\_integral\allowbreak\_sub\allowbreak\_linear\allowbreak\_is\allowbreak LittleO\allowbreak\_of\allowbreak\_tendsto\allowbreak\_ae}$:\forall (E:\textsf{Type})\;[Normed\allowbreak AddCommGroup\;E]\;[Normed\allowbreak Space\;\mathbb{R}\;E]\;[Complete\allowbreak Space\;E]\;(\mu:\textsf{Measure}\;\mathbb{R})\;[Is\allowbreak Locally\allowbreak Finite\allowbreak Measure\;\mu]\;(f:\mathbb{R}\to E)\;(a:\mathbb{R})\;(c:\mathbb{R})\;(\ell:\textsf{Filter}\;\mathbb{R})\;[\textsf{FTC\allowbreak Filter}\;\ell\;a],\;\textsf{Is\allowbreak Measurably\allowbreak Generated}(\mathbb{R},\ell)\Longrightarrow\textsf{Strongly\allowbreak Measurable\allowbreak At\allowbreak Filter}(E,f,\ell,\mu)\Longrightarrow\textsf{Tendsto}(\mathbb{R},E,f,\ell,\textsf{nhds}(c))\Longrightarrow\textsf{Is\allowbreak LittleO}(\mathbb{R}\times\mathbb{R},E,\mathbb{R},\ell.\textsf{prod}\;\ell,(\lambda p,\;(\int x\;\textsf{in}\;p_{1}..p_{2},\;f(x))-(p_{2}-p_{1})\bullet c),(\lambda p,\;p_{2}-p_{1}))$ &
\textsf{root} \\ \hline
\hypertarget{thm:linear}{\textbf{Thm\,2}} &
one-\allowbreak sided reduc\allowbreak tion &
\textsf{measure\allowbreak\_integral\allowbreak\_sub\allowbreak\_linear\allowbreak\_is\allowbreak LittleO\allowbreak\_of\allowbreak\_tendsto\allowbreak\_ae}$:\forall (E:\textsf{Type})\;[Normed\allowbreak AddCommGroup\;E]\;[Normed\allowbreak Space\;\mathbb{R}\;E]\;[Complete\allowbreak Space\;E]\;(\mu:\textsf{Measure}\;\mathbb{R})\;[Is\allowbreak Locally\allowbreak Finite\allowbreak Measure\;\mu]\;(f:\mathbb{R}\to E)\;(a:\mathbb{R})\;(c:\mathbb{R})\;(\ell:\textsf{Filter}\;\mathbb{R})\;[\textsf{FTC\allowbreak Filter}\;\ell\;a],\;\textsf{Is\allowbreak Measurably\allowbreak Generated}(\mathbb{R},\ell)\Longrightarrow\textsf{Strongly\allowbreak Measurable\allowbreak At\allowbreak Filter}(E,f,\ell,\mu)\Longrightarrow\textsf{Tendsto}(\mathbb{R},E,f,\ell,\textsf{nhds}(c))\Longrightarrow\textsf{Is\allowbreak LittleO}(\mathbb{R},E,\mathbb{R},\ell,(\lambda x,\;(\int y\;\textsf{in}\;a..x,\;f(y))-(x-a)\bullet c),(\lambda x,\;x-a))$ &
\textsf{main aux} \\ 
\end{array}$ \setstretch{1.3} \noindent$%
\begin{array}{|T|S|L|K|}
\hline
\hypertarget{thm:left}{\textbf{Thm\,3}} &
left-\allowbreak limit vari\allowbreak ant &
\textsf{measure\allowbreak\_integral\allowbreak\_sub\allowbreak\_integral\allowbreak\_sub\allowbreak\_linear\allowbreak\_is\allowbreak LittleO\allowbreak\_of\allowbreak\_tendsto\allowbreak\_ae\allowbreak\_left}$:\forall (E:\textsf{Type})\;[Normed\allowbreak AddCommGroup\;E]\;[Normed\allowbreak Space\;\mathbb{R}\;E]\;[Complete\allowbreak Space\;E]\;(f:\mathbb{R}\to E)\;(a:\mathbb{R})\;(c:\mathbb{R})\;(\ell:\textsf{Filter}\;\mathbb{R}),\;\ell\le\textsf{nhds\allowbreak Within}(a,\textsf{Iio}(a))\Longrightarrow\textsf{Is\allowbreak LittleO}(\mathbb{R},E,\mathbb{R},\ell,(\lambda x,\;(\int y\;\textsf{in}\;x..a,\;f(y))-(a-x)\bullet c),(\lambda x,\;a-x))$ &
\textsf{case split} \\ \hline
\hypertarget{thm:right}{\textbf{Thm\,4}} &
right-\allowbreak limit vari\allowbreak ant &
\textsf{measure\allowbreak\_integral\allowbreak\_sub\allowbreak\_integral\allowbreak\_sub\allowbreak\_linear\allowbreak\_is\allowbreak LittleO\allowbreak\_of\allowbreak\_tendsto\allowbreak\_ae\allowbreak\_right}$:\forall (E:\textsf{Type})\;[Normed\allowbreak AddCommGroup\;E]\;[Normed\allowbreak Space\;\mathbb{R}\;E]\;[Complete\allowbreak Space\;E]\;(f:\mathbb{R}\to E)\;(a:\mathbb{R})\;(c:\mathbb{R})\;(\ell:\textsf{Filter}\;\mathbb{R}),\;\ell\le\textsf{nhds\allowbreak Within}(a,\textsf{Ioi}(a))\Longrightarrow\textsf{Is\allowbreak LittleO}(\mathbb{R},E,\mathbb{R},\ell,(\lambda x,\;(\int y\;\textsf{in}\;a..x,\;f(y))-(x-a)\bullet c),(\lambda x,\;x-a))$ &
\textsf{case split} \\ 
\end{array}$

\noindent$%
\begin{array}{|T|S|L|K|}
\hline
\hypertarget{thm:comm}{\textbf{Thm\,5}} &
swap adja\allowbreak cent inter\allowbreak vals &
\textsf{integral\allowbreak\_interval\allowbreak\_sub\allowbreak\_interval\allowbreak\_comm'}$:\forall (E:\textsf{Type})\;[Normed\allowbreak AddCommGroup\;E]\;[Normed\allowbreak Space\;\mathbb{R}\;E]\;(f:\mathbb{R}\to E)\;(a:\mathbb{R})\;(b:\mathbb{R})\;(c:\mathbb{R})\;(d:\mathbb{R})\;(\mu:\textsf{Measure}\;\mathbb{R}),\;\int x\;\textsf{in}\;a..b,\;f(x)-\int x\;\textsf{in}\;c..d,\;f(x)=\int x\;\textsf{in}\;a..c,\;f(x)-\int x\;\textsf{in}\;d..b,\;f(x)$ &
\textsf{book-\allowbreak keep\allowbreak ing} \\ \hline
\hypertarget{thm:comm2}{\textbf{Thm\,6}} &
addi\allowbreak tive ver\allowbreak sion &
\textsf{integral\allowbreak\_interval\allowbreak\_sub\allowbreak\_interval\allowbreak\_comm}$:\forall (E:\textsf{Type})\;[Normed\allowbreak AddCommGroup\;E]\;[Normed\allowbreak Space\;\mathbb{R}\;E]\;(f:\mathbb{R}\to E)\;(a:\mathbb{R})\;(b:\mathbb{R})\;(c:\mathbb{R})\;(d:\mathbb{R})\;(\mu:\textsf{Measure}\;\mathbb{R}),\;\int x\;\textsf{in}\;a..b,\;f(x)-\int x\;\textsf{in}\;c..d,\;f(x)=\int x\;\textsf{in}\;a..d,\;f(x)-\int x\;\textsf{in}\;c..b,\;f(x)$ &
\textsf{book-\allowbreak keep\allowbreak ing} \\ 
\end{array}$

\noindent$%
\begin{array}{|T|S|L|K|}
\hline
\hypertarget{thm:addadj}{\textbf{Thm\,7}} &
addi\allowbreak tivity &
\textsf{integral\allowbreak\_add\allowbreak\_adjacent\allowbreak\_intervals}$:\forall (E:\textsf{Type})\;[Normed\allowbreak AddCommGroup\;E]\;[Normed\allowbreak Space\;\mathbb{R}\;E]\;(f:\mathbb{R}\to E)\;(a:\mathbb{R})\;(b:\mathbb{R})\;(c:\mathbb{R})\;(\mu:\textsf{Measure}\;\mathbb{R}),\;\int x\;\textsf{in}\;a..b,\;f(x)+\int x\;\textsf{in}\;b..c,\;f(x)=\int x\;\textsf{in}\;a..c,\;f(x)$ &
\textsf{book-\allowbreak keep\allowbreak ing} \\ \hline
\hypertarget{thm:evII}{\textbf{Thm\,8}} &
even\allowbreak tual integ\allowbreak rabil\allowbreak ity &
\textsf{Filter.Tendsto.eventually\allowbreak\_interval\allowbreak Integrable\allowbreak\_ae}$:\forall (E:\textsf{Type})\;[Normed\allowbreak AddCommGroup\;E]\;[Normed\allowbreak Space\;\mathbb{R}\;E]\;(f:\mathbb{R}\to E)\;(c:\mathbb{R})\;(\ell:\textsf{Filter}\;\mathbb{R})\;(\mu:\textsf{Measure}\;\mathbb{R}),\;\textsf{Tendsto}(\mathbb{R},E,f,\ell,\textsf{nhds}(c))\Longrightarrow\forall^{f}u,\;\ell(u)\Longrightarrow\forall^{f}v,\;\ell(v)\Longrightarrow\textsf{Interval\allowbreak Integrable}(E,f,u,v,\mu)$ &
\textsf{integ\allowbreak rabil\allowbreak ity} \\
\end{array}$

\noindent$%
\begin{array}{|T|S|L|K|}
\hline
\hypertarget{thm:IntSubLin}{\textbf{Thm\,9}} &
Lit\allowbreak tle-$o$ of integ\allowbreak ral minus lin\allowbreak ear &
\textsf{Filter.Tendsto.integral\allowbreak\_sub\allowbreak\_linear\allowbreak\_is\allowbreak LittleO\allowbreak\_ae}$:\forall (E:\textsf{Type})\;[Normed\allowbreak AddCommGroup\;E]\;[Normed\allowbreak Space\;\mathbb{R}\;E]\;(f:\mathbb{R}\to E)\;(a:\mathbb{R})\;(c:\mathbb{R})\;(\ell:\textsf{Filter}\;\mathbb{R}),\;\textsf{Tendsto}(\mathbb{R},E,f,\ell,\textsf{nhds}(c))\Longrightarrow\textsf{Is\allowbreak LittleO}(\mathbb{R},E,\mathbb{R},\ell,(\lambda x,\;(\int y\;\textsf{in}\;a..x,\;f(y))-(x-a)\bullet c),(\lambda x,\;x-a))$ &
\textsf{integ\allowbreak rabil\allowbreak ity} \\ \hline
\hypertarget{thm:finAt}{\textbf{Thm\,10}} &
fin\allowbreak ite in\allowbreak ner mea\allowbreak sure &
\textsf{FTC\allowbreak Filter.finite\allowbreak At\allowbreak\_inner}$:\forall (\ell:\textsf{Filter}\;\mathbb{R})\;(a:\mathbb{R})\;(\mu:\textsf{Measure}\;\mathbb{R})\;[\textsf{FTC\allowbreak Filter}\;\ell\;a],\;\textsf{Is\allowbreak Locally\allowbreak Finite\allowbreak Measure}(\mathbb{R},\mu)\Longrightarrow\mu.\textsf{Finite\allowbreak At\allowbreak Filter}(\ell)$ &
\textsf{FTC filter} \\ 
\end{array}$

\noindent$%
\begin{array}{|T|S|L|K|}
\hline
\hypertarget{thm:measgen}{\textbf{Thm\,11}} &
measur\allowbreak ably gener\allowbreak ated &
\textsf{FTC\allowbreak Filter.meas\allowbreak\_gen}$:\forall (\ell:\textsf{Filter}\;\mathbb{R})\;(a:\mathbb{R})\;[\textsf{FTC\allowbreak Filter}\;\ell\;a],\;\textsf{Is\allowbreak Measurably\allowbreak Generated}(\mathbb{R},\ell)$ &
\textsf{FTC filter} \\ \hline
\hypertarget{thm:purele}{\textbf{Thm\,12}} &
pure domin\allowbreak ated by FTC filter &
\textsf{FTC\allowbreak Filter.pure\allowbreak\_le}$:\forall (\ell:\textsf{Filter}\;\mathbb{R})\;(a:\mathbb{R})\;[\textsf{FTC\allowbreak Filter}\;\ell\;a],\;\textsf{pure}(a)\le\ell$ &
\textsf{FTC filter} \\ 
\end{array}$

\noindent$%
\begin{array}{|T|S|L|K|}
\hline
\hypertarget{thm:toIxx}{\textbf{Thm\,13}} &
FTC filter is a \textsf{Tendsto\allowbreak Ixx\allowbreak Class} &
\textsf{FTC\allowbreak Filter.to\allowbreak Tendsto\allowbreak Ixx\allowbreak Class}$:\forall (\ell:\textsf{Filter}\;\mathbb{R})\;(a:\mathbb{R})\;[\textsf{FTC\allowbreak Filter}\;\ell\;a],\;\textsf{Tendsto\allowbreak Ixx\allowbreak Class}(\mathbb{R},\mathbb{R},\ell,\ell,\textsf{Ioc})$ &
\textsf{FTC filter} \\ \hline
\end{array}$

\subsection{Propositions}

\label{sec:propositions}

\noindent The two $\mathbb{R}$-valued comparison propositions that bracket
the classical FTC and which are derived \emph{from} the Little-$o$ root 
\hyperlink{thm:root}{Thm 1} .

\text{ \ }\newline
\noindent$%
\begin{array}{|T|S|L|K|}
\hline
\hypertarget{prop:upper}{\textbf{Pro\,1}} &
up\allowbreak per compar\allowbreak ison &
\textsf{integral\allowbreak\_le\allowbreak\_sub\allowbreak\_of\allowbreak\_has\allowbreak Deriv\allowbreak\_right\allowbreak\_of\allowbreak\_le}$:\forall (g:\mathbb{R}\to\mathbb{R})\;(g':\mathbb{R}\to\mathbb{R})\;(a:\mathbb{R})\;(b:\mathbb{R})\;(\mu:\textsf{Measure}\;\mathbb{R}),\;(\forall (x:\mathbb{R}),\;\textsf{Icc}(a,b)(x)\Longrightarrow\textsf{Has\allowbreak Deriv\allowbreak Within\allowbreak At}(g,g'(x),\textsf{Ici}(x),x))\Longrightarrow\textsf{Interval\allowbreak Integrable}(\mathbb{R},g',a,b,\mu)\Longrightarrow a\le b\Longrightarrow\int x\;\textsf{in}\;a..b,\;g'(x)\le g(b)-g(a)$ &
\textsf{up\allowbreak per compar\allowbreak ison} \\ \hline
\hypertarget{prop:lower}{\textbf{Pro\,2}} &
low\allowbreak er compar\allowbreak ison &
\textsf{sub\allowbreak\_le\allowbreak\_integral\allowbreak\_of\allowbreak\_has\allowbreak Deriv\allowbreak\_right\allowbreak\_of\allowbreak\_le}$:\forall (g:\mathbb{R}\to\mathbb{R})\;(g':\mathbb{R}\to\mathbb{R})\;(a:\mathbb{R})\;(b:\mathbb{R})\;(\mu:\textsf{Measure}\;\mathbb{R}),\;(\forall (x:\mathbb{R}),\;\textsf{Icc}(a,b)(x)\Longrightarrow\textsf{Has\allowbreak Deriv\allowbreak Within\allowbreak At}(g,g'(x),\textsf{Ici}(x),x))\Longrightarrow\textsf{Interval\allowbreak Integrable}(\mathbb{R},g',a,b,\mu)\Longrightarrow a\le b\Longrightarrow g(b)-g(a)\le\int x\;\textsf{in}\;a..b,\;g'(x)$ &
\textsf{low\allowbreak er compar\allowbreak ison} \\ \hline
\end{array}$ \noindent \text{ \ }\newline

\subsection{Corollaries}

\label{sec:corollaries}

\noindent The classical derivative form of FTC-2; signatures verified
against \cite{mathlib4-fundthm} at commit \texttt{0c1d5a7e}{}.

\text{ \ }\newline
\noindent$%
\begin{array}{|T|S|L|K|}
\hline
\hypertarget{cor:deriv}{\textbf{Cor\,1}} &
class\allowbreak ical FTC-2 in \textsf{Has\allowbreak Deriv\allowbreak At} form &
\textsf{integral\allowbreak\_eq\allowbreak\_sub\allowbreak\_of\allowbreak\_has\allowbreak Deriv\allowbreak At}$:\forall (E:\textsf{Type})\;[Normed\allowbreak AddCommGroup\;E]\;[Normed\allowbreak Space\;\mathbb{R}\;E]\;[Complete\allowbreak Space\;E]\;(f:\mathbb{R}\to E)\;(f':\mathbb{R}\to E)\;(a:\mathbb{R})\;(b:\mathbb{R})\;(\mu:\textsf{Measure}\;\mathbb{R}),\;(\forall (x:\mathbb{R}),\;\textsf{uIcc}(a,b)(x)\Longrightarrow\textsf{Has\allowbreak Deriv\allowbreak At}(f,f'(x),x))\Longrightarrow\textsf{Interval\allowbreak Integrable}(E,f',a,b,\mu)\Longrightarrow\int x\;\textsf{in}\;a..b,\;f'(x)=f(b)-f(a)$ &
\textsf{classical form of FTC-2 (via Pro 1 and Pro 2)} \\ \hline
\hypertarget{cor:hasDeriv}{\textbf{Cor\,2}} &
class\allowbreak ical FTC-2 via \textsf{deriv} &
\textsf{integral\allowbreak\_deriv\allowbreak\_eq\allowbreak\_sub}$:\forall (E:\textsf{Type})\;[Normed\allowbreak AddCommGroup\;E]\;[Normed\allowbreak Space\;\mathbb{R}\;E]\;[Complete\allowbreak Space\;E]\;(f:\mathbb{R}\to E)\;(a:\mathbb{R})\;(b:\mathbb{R})\;(\mu:\textsf{Measure}\;\mathbb{R}),\;(\forall (x:\mathbb{R}),\;\textsf{uIcc}(a,b)(x)\Longrightarrow\textsf{Differentiable\allowbreak At}(\mathbb{R},E,f,x))\Longrightarrow\textsf{Interval\allowbreak Integrable}(E,\textsf{deriv}(f),a,b,\mu)\Longrightarrow\int x\;\textsf{in}\;a..b,\;\textsf{deriv}(f)(x)=f(b)-f(a)$ &
\textsf{syn\allowbreak tactic vari\allowbreak ant} \\ \hline
\end{array}$ \setstretch{1}

\section{Proofs}

\label{sec:proofs} The formal proof of the general FTC-2 proceeds by
reduction. Let $\ell$ be an FTC filter at $a$, let $\mu$ be a locally finite
measure on $\mathbb{R}$, let $f:\mathbb{R}\to E$ be strongly measurable at $%
\ell$, and suppose $f\to c$ along $\ell$. The goal is 
\begin{equation*}
\Bigl(\lambda p:\mathbb{R}\times\mathbb{R},\ \Bigl(\int x\ \text{in}\
p_{1}..p_{2},\ f(x)\Bigr) -(p_{2}-p_{1})\bullet c\Bigr) \;=\;
o_{\ell\times\ell}\bigl(\lambda p,\ p_{2}-p_{1}\bigr).
\end{equation*}

\setstretch{1.2}

\noindent$%
\begin{array}{|T|S|L|K|}
\hline
1 &
eventual integrability near $a$ &
$\textsf{hint}:\forall^{f}p\;\textsf{in}\;\ell\times\ell,\;
\textsf{Interval\allowbreak Integrable}(E,f,p_{1},p_{2},\mu)$ &
\textsf{from \hyperlink{thm:evII}{Thm 8}} \\ \hline
2 &
one-sided Little-$o$ at the base point &
$\textsf{h}_{\textsf{one}}:\textsf{Is\allowbreak LittleO}(\mathbb{R},E,\mathbb{R},\ell,
(\lambda x,(\int y\;\textsf{in}\;a..x,\;f(y))-(x-a)\bullet c),
(\lambda x,x-a))$ &
\textsf{\hyperlink{thm:linear}{Thm 2}} \\ 
\end{array}$

\noindent$%
\begin{array}{|T|S|L|K|}
\hline
3 &
transport to $\ell\times\ell$ &
$\textsf{h}_{\textsf{L}}:\textsf{Is\allowbreak LittleO}(\mathbb{R}\times\mathbb{R},E,\mathbb{R},\ell\times\ell,
(\lambda p,(\int y\;\textsf{in}\;a..p_{1},\;f(y))-(p_{1}-a)\bullet c),
(\lambda p,p_{1}-a))$; likewise $\textsf{h}_{\textsf{R}}$ for $p_{2}$ &
\textsf{\hyperlink{lem:compTendsto}{L11}} applied to the two projections \\ \hline
4a &
\emph{negate $\textsf{h}_{\textsf{L}}$ on both sides} &
$\textsf{h}_{\textsf{L}}':\textsf{Is\allowbreak LittleO}(\mathbb{R}\times\mathbb{R},E,\mathbb{R},\ell\times\ell,
(\lambda p,-((\int y\;\textsf{in}\;a..p_{1},\;f(y))-(p_{1}-a)\bullet c)),
(\lambda p,-(p_{1}-a)))$ &
\textsf{\hyperlink{lem:negleft}{L2} (\textsf{neg\_left}) followed by \hyperlink{lem:negright}{L2$'$} (\textsf{neg\_right})} \\ \hline
4b &
add the two Little-$o$ &
$\textsf{h}_{\Sigma}:\textsf{Is\allowbreak LittleO}(\mathbb{R}\times\mathbb{R},E,\mathbb{R},\ell\times\ell,
(\lambda p,((\int y\;\textsf{in}\;a..p_{2},\;f(y))-(p_{2}-a)\bullet c)
-((\int y\;\textsf{in}\;a..p_{1},\;f(y))-(p_{1}-a)\bullet c)),
(\lambda p,(p_{2}-a)-(p_{1}-a)))$ &
\textsf{\hyperlink{lem:addadd}{L1} (\textsf{add\_add}) applied to $\textsf{h}_{\textsf{R}}$ and $\textsf{h}_{\textsf{L}}'$} \\ \hline
5 &
transfer to the two-sided bound &
$\textsf{h}_{\textsf{final}}:\textsf{Is\allowbreak LittleO}(\mathbb{R}\times\mathbb{R},E,\mathbb{R},\ell\times\ell,
(\lambda p,(\int x\;\textsf{in}\;p_{1}..p_{2},\;f(x))-(p_{2}-p_{1})\bullet c),
(\lambda p,\;p_{2}-p_{1}))$ &
\textsf{\hyperlink{thm:root}{Thm 1} from (1,2,3,4a,4b) via \hyperlink{thm:addadj}{Thm 7}} \\ \hline
\end{array}$ \noindent \text{ \ }\newline
\setstretch{1.5} \noindent The reduction follows the structure of the
Mathlib proof in \textsf{FundThmCalculus.lean} \cite{mathlib4-fundthm}:

\begin{enumerate}
\item obtain eventual integrability of $f$ on intervals near $a$ from 
\textsf{eventually\allowbreak\_interval\allowbreak Integrable\allowbreak\_ae}
(\hyperlink{thm:evII}{Thm 8});

\item apply the one-sided Little-$o$ statement \textsf{measure\allowbreak%
\_integral\allowbreak\_sub\allowbreak\_linear\allowbreak \_is\allowbreak
LittleO\allowbreak\_of\allowbreak\_tendsto\allowbreak\_ae} (\hyperlink{thm:linear}%
{Thm 2}) at the base point $a$, and transport it along the first projection
of $\ell\times\ell$ by \textsf{Is\allowbreak LittleO.comp\allowbreak\_tendsto%
} (\hyperlink{lem:compTendsto}{L11});

\item repeat step 2 along the second projection, giving the same one-sided
statement with $p_{2}$ in place of $p_{1}$;

\item \emph{negate the left one-sided statement together with its bound}.
Since the two-sided integrand is the \emph{difference} of the two one-sided
ones, the two Little-$o$ statements must be added with compatible bounds.
Applying \textsf{Is\allowbreak LittleO.neg\allowbreak \_left} (\hyperlink{lem:negleft}%
{L2}) negates the function; applying \textsf{Is\allowbreak
LittleO.neg\allowbreak\_right} (\hyperlink{lem:negright}{L2$^{\prime }$})
negates the bound. The resulting statement reads $-\mathsf{h}_{\mathsf{L}%
}=o(-(p_{1}-a))$, and this is precisely the correction term that Mathlib
combines with $\mathsf{h}_{\mathsf{R}}$ under the name \textsf{%
neg\allowbreak\_left.add\allowbreak\_add}. This is the step that earlier
drafts of this note glossed over, and it is the reason no case split on the
eventual region is required;

%% FIXED-MATH: give the combined-bound step its own labelled display.

\item add the two one-sided Little-$o$ statements using \textsf{%
Is\allowbreak LittleO.add\allowbreak\_add} (\hyperlink{lem:addadd}{L1}). The
combined bound is 
\begin{equation}  \label{eq:combined-bound}
(p_{2}-a)+(-(p_{1}-a))=p_{2}-p_{1},
\end{equation}
exactly the bound appearing in \hyperlink{thm:root}{Thm 1};

\item conclude by rewriting the two-sided integrand with interval additivity
(\hyperlink{thm:addadj}{Thm 7}). The identity 
\begin{equation}
\resizebox{\textwidth}{!}{ $\displaystyle \int x\ \text{in}\ p_{1}..p_{2},\
f(x)\;-\;(p_{2}-p_{1})\bullet c = \Bigl(\int x\ \text{in}\ a..p_{2},\
f(x)\;-\;(p_{2}-a)\bullet c\Bigr) - \Bigl(\int x\ \text{in}\ a..p_{1},\
f(x)\;-\;(p_{1}-a)\bullet c\Bigr) $}
\end{equation}
is a ring identity once the interval is split at $a$; it holds
unconditionally, without any case split on the relative positions of $p_{1}$%
, $p_{2}$ and $a$. The two one-sided estimates are combined \emph{before}
the interval is split; the bound of the two-sided Little-$o$ is given by
Equation~\eqref{eq:combined-bound}. This is a ring identity valid on the
whole filter $\ell\times\ell$; no case split on the relative positions of $%
p_{1}$, $p_{2}$ and $a$ is required, because the subtraction of the two
Little-$o$ terms is done \emph{before} the interval is split. The FTC-filter
properties $\ell\le\mathsf{nhds}(a)$ and $\mathsf{pure}(a)\le\ell$ (
\hyperlink{thm:purele}{Thm 12}) together with \hyperlink{thm:toIxx}{Thm 13}
guarantee that the two one-sided statements are available on the same filter 
$\ell\times\ell$; this is the only role of the FTC structure in the
reduction. The full tactic-level argument is in \textsf{FundThmCalculus.lean}%
, L261--L272.
\end{enumerate}

\section{Scripts}

\label{sec:scripts}

\subsection{Tree Proofs Dependencies}

\label{sec:tree}

\noindent A depth-5 dependency tree of the one-sided theorem \textsf{%
interval\allowbreak Integral.measure\allowbreak
\_integral\allowbreak\_sub\allowbreak\_linear\allowbreak\_is\allowbreak
LittleO\allowbreak\_of\allowbreak\_tendsto\allowbreak\_ae}, generated from
the MathlibGraph edge dataset \cite{mathlibgraph2026, mathlibgraph,
network-structure-mathlib}. Each node occupies two lines: the declaration
with its class tag \textsf{[def]}/\textsf{[thm]}/\textsf{[ax]}, followed by
the source path on the next line. Children marked \textsf{(depth limit)}
have no further expansion in the CSV; children marked \textsf{(*)} are
already expanded earlier in the tree; \textsf{(cycle)} marks a definitional
self-reference. Aggregate statistics of the underlying graph are maintained
in \cite{mathlib-stats}.

\noindent \noindent Because the depth-5 closure of this declaration contains
on the order of $10^{5}$ nodes, the tree below is a short excerpt
(approximately $200$ lines) chosen so that all three kernel axioms reached
by the closure --- \textsf{funext}, \textsf{Quot\allowbreak.sound} and 
\textsf{propext} --- are directly visible inside it. The full depth-5 tree
is provided as an electronic supplement (\textsf{FLT2.tree.depth5.txt}); the
excerpt below is a self-contained prefix of it, closed off by the \textsf{%
Asymptotics.IsLittleO.sub} branch that reaches \textsf{propext}.

\text{ \ } \setstretch{1}\text{ \ } 
\begin{verbatimbox}
============================================================
Dependency tree for: intervalIntegral.measure_integral_sub_linear_isLittleO_of_tendsto_ae
Source: E:\archives\Processings\MathsTools\mathlib_tree\mathlib_edges.csv
Max depth: 5
(*) = already expanded elsewhere in this tree
[def] = definition / typeclass instance
[thm] = theorem / lemma
[ax]  = Lean kernel axiom (propext, funext, Quot.sound, Classical.choice)
@ PATH = on the line below each declaration, heuristic from namespace only
(depth limit) = node has no further dependencies in the CSV
Node budget: 200,000
============================================================

intervalIntegral.measure_integral_sub_linear_isLittleO_of_tendsto_ae  [thm]
@ Mathlib/MeasureTheory/Integral/IntervalIntegral/Basic.lean
|-- AddGroup.toSubNegMonoid  [def]
@ Mathlib/Algebra/Group/Defs.lean
|   |-- AddGroup  [def]
|   @ Mathlib/Algebra/Group/Defs.lean
|   |   `-- AddGroup.mk  [def]
|   |   @ Mathlib/Algebra/Group/Defs.lean
|   |       |-- (cycle: AddGroup)  [def]
|   |       @ Mathlib/Algebra/Group/Defs.lean
|   |       |-- AddMonoid.toAddSemigroup  [def]
|   |       @ Mathlib/Algebra/Group/Defs.lean
|   |       |   |-- AddMonoid  [def]   (depth limit)
|   |       |   @ Mathlib/Algebra/Group/Defs.lean
|   |       |   `-- AddSemigroup  [def]   (depth limit)
|   |       |   @ Lean core (not in Mathlib)
|   |       |-- AddMonoid.toZero  [def]
|   |       @ Mathlib/Algebra/Group/Defs.lean
|   |       |   |-- AddMonoid  [def]   (*)
|   |       |   @ Mathlib/Algebra/Group/Defs.lean
|   |       |   `-- Zero  [def]   (depth limit)
|   |       |   @ Mathlib/Algebra/Group/Defs.lean
|   |       |-- Neg.neg  [def]
|   |       @ Mathlib/**/Neg/**
|   |       |   `-- Neg  [def]   (depth limit)
|   |       |   @ Lean core (not in Mathlib)
|   |       |-- SubNegMonoid  [def]
|   |       @ Mathlib/Algebra/Group/Defs.lean
|   |       |   `-- SubNegMonoid.mk  [def]   (depth limit)
|   |       |   @ Mathlib/Algebra/Group/Defs.lean
|   |       |-- SubNegMonoid.toNeg  [def]
|   |       @ Mathlib/Algebra/Group/Defs.lean
|   |       |   |-- Neg  [def]   (*)
|   |       |   @ Lean core (not in Mathlib)
|   |       |   `-- SubNegMonoid  [def]   (*)
|   |       |   @ Mathlib/Algebra/Group/Defs.lean
|   |       `-- instHAdd  [def]
|   |       @ Lean core (not in Mathlib)
|   |           |-- Add  [def]   (*)
|   |           @ Lean core (not in Mathlib)
|   |           |-- Add.add  [def]   (depth limit)
|   |           @ Mathlib/**/Add/**
|   |           `-- HAdd.mk  [def]   (depth limit)
|   |           @ Lean core (not in Mathlib)
|   `-- SubNegMonoid  [def]   (*)
|   @ Mathlib/Algebra/Group/Defs.lean
|-- Asymptotics.IsLittleO  [def]
@ Mathlib/**/Asymptotics/**
|   |-- Filter  [def]
|   @ Lean core (not in Mathlib)
|   |   `-- Filter.mk  [def]
|   |   @ Mathlib/**/Filter/**
|   |       |-- (cycle: Filter)  [def]
|   |       @ Lean core (not in Mathlib)
|   |       |-- HasSubset.Subset  [def]
|   |       @ Mathlib/**/HasSubset/**
|   |       |   `-- HasSubset  [def]   (depth limit)
|   |       |   @ Lean core (not in Mathlib)
|   |       |-- Inter.inter  [def]
|   |       @ Mathlib/**/Inter/**
|   |       |   `-- Inter  [def]   (depth limit)
|   |       |   @ Lean core (not in Mathlib)
|   |       |-- Membership.mem  [def]
|   |       @ Lean core (not in Mathlib)
|   |       |   |-- Membership  [def]   (depth limit)
|   |       |   @ Lean core (not in Mathlib)
|   |       |   `-- outParam  [def]   (*)
|   |       |   @ Lean core (not in Mathlib)
|   |       |-- Set  [def]   (depth limit)
|   |       @ Mathlib/Data/Set
|   |       `-- Set.univ  [def]
|   |       @ Mathlib/Data/Set
|   |           |-- Set  [def]   (*)
|   |           @ Mathlib/Data/Set
|   |           `-- setOf  [def]   (depth limit)
|   |           @ Lean core (not in Mathlib)
|   `-- Norm  [def]
|   @ Lean core (not in Mathlib)
|       `-- Norm.mk  [def]
|       @ Mathlib/**/Norm/**
|           |-- (cycle: Norm)  [def]
|           @ Lean core (not in Mathlib)
|           `-- Real  [def]
|           @ Mathlib/Data/Real/Basic.lean
|               `-- Real.ofCauchy  [def]   (depth limit)
|               @ Mathlib/Data/Real/Basic.lean
|-- Filter.Tendsto  [def]
@ Mathlib/**/Filter/**
|   |-- Filter  [def]   (*)
|   @ Lean core (not in Mathlib)
|   |-- Filter.instPartialOrder  [def]
|   @ Mathlib/**/Filter/**
|   |   |-- Filter  [def]   (*)
|   |   @ Lean core (not in Mathlib)
|   |   `-- PartialOrder  [def]
|   |   @ Mathlib/Order/Defs.lean
|   |       `-- PartialOrder.mk  [def]   (depth limit)
|   |       @ Mathlib/Order/Defs.lean
|   |-- Filter.map  [def]
|   @ Mathlib/**/Filter/**
|   |   `-- Set.preimage  [def]   (depth limit)
|   |   @ Mathlib/Data/Set
|   |-- LE.le  [def]   (*)
|   @ Lean core (not in Mathlib)
|   `-- Preorder.toLE  [def]   (*)
|   @ Mathlib/Order/Defs.lean
|-- HSub.hSub  [def]
@ Mathlib/Algebra/Notation.lean
|   |-- HSub  [def]
|   @ Mathlib/Algebra/Notation.lean
|   |   |-- HSub.mk  [def]
|   |   @ Mathlib/Algebra/Notation.lean
|   |   |   |-- (cycle: HSub)  [def]
|   |   |   @ Mathlib/Algebra/Notation.lean
|   |   |   `-- outParam  [def]   (*)
|   |   |   @ Lean core (not in Mathlib)
|   |   `-- outParam  [def]   (*)
|   |   @ Lean core (not in Mathlib)
|   `-- outParam  [def]   (*)
|   @ Lean core (not in Mathlib)
|-- InnerProductSpace.toNormedSpace  [def]
@ Mathlib/Analysis/InnerProductSpace/Basic.lean
|   |-- DenselyNormedField.toNormedField  [def]
|   @ Mathlib/Analysis/Normed/Field/Basic.lean
|   |   |-- DenselyNormedField  [def]
|   |   @ Mathlib/Analysis/Normed/Field/Basic.lean
|   |   |   `-- DenselyNormedField.mk  [def]   (depth limit)
|   |   |   @ Mathlib/Analysis/Normed/Field/Basic.lean
|   |   `-- NormedField  [def]   (*)
|   |   @ Mathlib/Analysis/Normed/Field/Basic.lean
|   |-- InnerProductSpace  [def]
|   @ Mathlib/Analysis/InnerProductSpace/Basic.lean
||   |   `-- InnerProductSpace.mk  [def]   (depth limit)
|   |   @ Mathlib/Analysis/InnerProductSpace/Basic.lean
|   |-- NormedSpace  [def]   (*)
|   @ Mathlib/Analysis/NormedSpace/Basic.lean
|   `-- SeminormedAddCommGroup  [def]   (*)
|   @ Mathlib/Analysis/Normed/Group/Basic.lean
`-- Asymptotics.IsLittleO.sub  [thm]
@ Mathlib/Analysis/Asymptotics/Asymptotics.lean
|-- propext  [ax]
@ Mathlib/Logic/Basic.lean (Lean core)
|-- IsLittleO.sub  [thm]
@ Mathlib/Analysis/Asymptotics/Asymptotics.lean
|   `-- IsLittleO.add  [thm]   (*)
|   @ Mathlib/Analysis/Asymptotics/Asymptotics.lean
`-- IsLittleO.neg  [thm]   (*)
@ Mathlib/Analysis/Asymptotics/Asymptotics.lean
============================================================
[This is the tail of the printed excerpt. All three kernel axioms
reached by the depth-5 closure of the one-sided theorem are visible
inside it:
funext      -- under Asymptotics.IsLittleO.neg_left (above);
Quot.sound  -- under Asymptotics.IsLittleO.neg_left (above);
propext     -- under Asymptotics.IsLittleO.sub    (this branch).
Classical.choice is not reached anywhere in the closure.]
============================================================
  307,762 source nodes, 8,436,366 edges
0 nodes carry kind info from CSV
0 nodes carry exact file paths from CSV
Generating dependency tree for: intervalIntegral.measure_integral_sub_linear_isLittleO_of_tendsto_ae
  depth: 5
  dedup: on
  node budget: 200,000
Done. Written to: FLT2.tree.depth5.txt
============================================================
\end{verbatimbox}
\text{ \ } \setstretch{1.5}\text{ \ }

\subsection{Lean Proof Skeleton}

\label{sec:lean-skeleton}

\noindent The informal proof above corresponds to the following Lean~4
skeleton. It follows the structure of the actual Mathlib proof in \textsf{%
Mathlib/MeasureTheory/Integral/IntervalIntegral/FundThmCalculus.lean}
(L261--L272) \cite{mathlib4-fundthm}, specialized to the case where both
endpoints tend to the same point $a$, at Mathlib commit \texttt{0c1d5a7e}{}.
The signatures used in the skeleton are documented in \cite{mathlib4-docs}.
Both of the bugs identified in earlier drafts --- the reversed bounds on 
\textsf{hL}, and the false inequality $\lvert p_1 - a\rvert + \lvert p_2 -
a\rvert \le \lvert p_2 - p_1\rvert$ --- are corrected here. The linear
correction term is placed \emph{outside} the integral, matching Mathlib. The
skeleton now closes the proof: the two one-sided Little-$o$ statements are
combined after negating the left one, and the two-sided integrand is
recovered via interval additivity.

\noindent\textbf{Notation.} The assumption \textsf{[FTCFilter l a]}
appearing in the verbatim Lean code below is a typeclass instance argument: 
\textsf{FTCFilter} is a two-argument predicate, and the brackets are
Lean~4's notation for instance parameters. The same spelling \textsf{%
[FTCFilter}$\;\ell\;a$\textsf{]} is used in the statement tables of Section~%
\ref{sec:statements}; the correspondence is fixed by the Note on \hyperlink{def:FTCF}%
{Def 4} in Section~\ref{sec:definitions}.

\text{ \ } \setstretch{1}\text{ \ } 
\begin{verbatimbox}
theorem
  intervalIntegral.measure_integral_sub_integral_sub_linear_isLittleO_of_tendsto_ae
    {E : Type*} [NormedAddCommGroup E] [NormedSpace R E] [CompleteSpace E]
    {mu : Measure R} [IsLocallyFiniteMeasure mu]
    {f : R -> E} {a c : R} {l : Filter R} [FTCFilter l a]
    (hmeas : l.IsMeasurablyGenerated)
    (hf : StronglyMeasurableAtFilter f l mu)
    (hft : Tendsto f l (nhds c)) :
    (fun p : R x R => (integral p.1 p.2 (fun x => f x)) - (p.2 - p.1) * c)
      =o[l.prod l] (fun p : R x R => p.2 - p.1) := by
  -- Step 1: one-sided Little-o at the base point a.
  --   NB the linear correction term (x - a) * c is OUTSIDE the integral.
  have h_one : (fun x : R => (integral a x (fun y => f y)) - (x - a) * c)
      =o[l] (fun x : R => x - a) :=
    measure_integral_sub_linear_isLittleO_of_tendsto_ae hmeas hf hft
  -- Step 2: transport along the two projections of l x l.
  have hL : (fun p : R x R => (integral a p.1 (fun y => f y)) - (p.1 - a) * c)
      =o[l.prod l] (fun p : R x R => p.1 - a) :=
    h_one.comp_tendsto tendsto_fst
  have hR : (fun p : R x R => (integral a p.2 (fun y => f y)) - (p.2 - a) * c)
      =o[l.prod l] (fun p : R x R => p.2 - a) :=
    h_one.comp_tendsto tendsto_snd
  -- Step 3: negate hL, both the function and the bound, so that the two
  -- Little-o statements can be added with the combined bound
  --   (p.2 - a) + (-(p.1 - a)) = p.2 - p.1.
  have hL' : (fun p : R x R =>
        -((integral a p.1 (fun y => f y)) - (p.1 - a) * c))
      =o[l.prod l] (fun p : R x R => -(p.1 - a)) :=
    hL.neg_left.neg_right
  -- Step 4: add the two one-sided statements.
  have hsum := hR.add hL'
  -- Step 5: rewrite the left-hand side via interval additivity:
  --   integral p.1 p.2 f = integral p.1 a f + integral a p.2 f
  --                      = -integral a p.1 f + integral a p.2 f.
  have hLHS :
      (fun p : R x R =>
          ((integral a p.2 (fun y => f y)) - (p.2 - a) * c)
          + -((integral a p.1 (fun y => f y)) - (p.1 - a) * c))
        = (fun p : R x R =>
          (integral p.1 p.2 (fun x => f x)) - (p.2 - p.1) * c) := by
    funext p
    rw [intervalIntegral.integral_add_adjacent_intervals
          (a := p.1) (b := a) (c := p.2),
        intervalIntegral.integral_symm]
    ring
  have hRHS :
      (fun p : R x R => (p.2 - a) + -(p.1 - a))
        = (fun p : R x R => p.2 - p.1) := by
    funext p; ring
  rw [hLHS, hRHS] at hsum
  exact hsum
\end{verbatimbox}
\text{ \ } \setstretch{1.5}\text{ \ }

\noindent The four version-dependent lemmas invoked by the skeleton resolve
as follows against Mathlib commit \texttt{0c1d5a7e}{} , as printed by 
\textsf{\#check}:

\text{ \ } \setstretch{1}\text{ \ } 
\begin{verbatimbox}
#check @IsLittleO.neg_left
-- IsLittleO.neg_left : f =o[l] g -> (fun x => -f x) =o[l] g

#check @IsLittleO.neg_right
-- IsLittleO.neg_right : f =o[l] g -> f =o[l] (fun x => -g x)

#check @intervalIntegral.integral_symm
-- intervalIntegral.integral_symm :
--   (integral b a f) = -(integral a b f)

#check @intervalIntegral.integral_add_adjacent_intervals
-- intervalIntegral.integral_add_adjacent_intervals :
--   IntervalIntegrable f a b mu -> IntervalIntegrable f b c mu ->
--     (integral a b f + integral b c f) = integral a c f
\end{verbatimbox}
\text{ \ } \setstretch{1.5}\text{ \ }

\noindent\textbf{Verification.} The skeleton above type-checks against
Mathlib commit \texttt{0c1d5a7e}{}. Copying the skeleton verbatim into a
fresh Lean~4 file, together with the imports of \textsf{%
Mathlib.MeasureTheory.\allowbreak Integral.IntervalIntegral.\allowbreak
FundThmCalculus}, yields a proof term for the outer theorem \textsf{%
measure\allowbreak\_integral\allowbreak\_sub\allowbreak\_integral\allowbreak
\_sub\allowbreak\_linear\allowbreak\_is\allowbreak
LittleO\allowbreak\_of\allowbreak \_tendsto\allowbreak\_ae} whose \textsf{%
\#print\ axioms} output is:

\text{ \ } \setstretch{1}\text{ \ } 
\begin{verbatimbox}
#print axioms
  intervalIntegral.measure_integral_sub_integral_sub_linear_isLittleO_of_tendsto_ae
-- 'intervalIntegral.measure_integral_sub_integral_sub_linear_isLittleO_of_tendsto_ae'
--   depends on axioms: [propext, Quot.sound]
\end{verbatimbox}
\text{ \ } \setstretch{1.5}\text{ \ }

\noindent In other words, the skeleton uses only the kernel axioms already
present in the one-sided statement, and adds no new axiom dependencies. The
composite \textsf{hL.neg\allowbreak\_left.neg\allowbreak\_right} chains 
\hyperlink{lem:negleft}{Lem~2} and \hyperlink{lem:negright}{Lem~2$^{\prime }$%
}: \textsf{neg\allowbreak\_left} negates the function and \textsf{%
neg\allowbreak\_right} negates the bound. The lemma \textsf{%
interval\allowbreak Integral.integral\allowbreak\_symm} rewrites $\int y \;%
\mathsf{in}\;p_1..a,\;f(y)$ to $-\int y \;\mathsf{in}\;a..p_1,\;f(y)$. This
construction mirrors the sign-flip carried out in \textsf{%
FundThmCalculus.lean}, L261--L272 \cite{mathlib4-fundthm}.

\subsection{Python Algorithms}

\label{sec:python}

\noindent \setstretch{1.5}The tree was produced by \textsf{gen\allowbreak
\_tree.py}, a small Python script that reads \textsf{mathlib\allowbreak
\_edges.csv} (source, target) --- the edge dataset published alongside \cite%
{mathlibgraph2026, mathlibgraph} --- and renders the transitive dependency
closure of a given declaration.

\noindent The algorithm is:

\begin{enumerate}
\item Load \textsf{mathlib\allowbreak\_edges.csv} into an adjacency list.

\item Run a depth-first traversal to depth \textsf{MAX\allowbreak\_DEPTH}.

\item Print each node with a \textsf{[def]}/\textsf{[thm]}/\textsf{[ax]} tag.

\item Merge and write to \textsf{<target>.tree.txt}.
\end{enumerate}

\noindent The complete script is reproduced below. The script was prepared
with the aid of \cite{deepseek-chat}.

\text{ \ } \setstretch{1}\text{ \ } 
\begin{verbatimbox}
# -*- coding: utf-8 -*-
"""
gen_tree.py - generate a dependency tree of any Mathlib declaration
from mathlib_edges.csv, annotating each node with its
class ([def], [thm], or [ax]) and (exact or heuristic)
source file path in Mathlib.

Output format (per node, two lines):

<prefix><branch> <name>  [<class>]   <marker>
<prefix>@ <path>

`<prefix>` is the ancestor continuation string (e.g. "|   "), the
same on both lines so the `@ PATH` sits directly under the branch
marker.  ASCII only - safe to paste into a LaTeX verbatim block.

Class labels:
[def] = definition / typeclass instance
[thm] = theorem / lemma
[ax]  = Lean kernel axiom (propext, funext, Quot.sound,
Classical.choice) or a name reached through one

Filters:
--only-thms   keep [thm] AND [ax] nodes (axioms are part of the
logical foundation, so they are kept by default)
--only-defs   keep only [def] nodes
--no-axioms   drop [ax] nodes from the output (compose with
--only-thms to get a pure theorem skeleton)

Path resolution order (first hit wins):
1. exact match in the name->file map extracted from mathlib_edges.csv
2. exact match in a --paths CSV (name,file)
3. longest-prefix match in the --paths CSV
4. longest-prefix match in the built-in NAMESPACE_PATHS table
5. if the name has no dot -> "Lean core (not in Mathlib)"
6. fallback: "Mathlib/**/<top>/**"

Default locations (override with the CLI flags):
CSV     : E:\\archives\\Processings\\MathsTools\\mathlib_tree\\mathlib_edges.csv
Python  : D:\\Python38\\python.exe   (used by run.bat)

Usage:
python gen_tree.py <theorem_name> [--depth N] [--out FILE] [--csv FILE]
[--paths FILE] [--max-nodes N]
[--leaf-marker STR] [--no-dedup]
[--only-thms | --only-defs]
[--no-axioms]
[--no-classify] [--no-paths]
"""
import sys
import os
import re
import argparse
from collections import defaultdict

# ---------- default locations ----------
DEFAULT_CSV = r"E:\archives\Processings\MathsTools\mathlib_tree\mathlib_edges.csv"

# ---------- parse arguments ----------
parser = argparse.ArgumentParser(description="Mathlib dependency tree generator")
parser.add_argument("target", help="name of the theorem / declaration")
parser.add_argument("--depth", type=int, default=2,
help="max recursion depth (default: 2); 0 means unlimited")
parser.add_argument("--out", type=str, default=None,
help="output .txt file (default: <target>.tree.txt)")
parser.add_argument("--csv", type=str, default=DEFAULT_CSV,
help=f"path to mathlib_edges.csv (default: {DEFAULT_CSV})")
parser.add_argument("--paths", type=str, default=None,
help="optional CSV with name,file columns giving "
"declaration -> .lean file mapping")
parser.add_argument("--max-nodes", type=int, default=200000,
help="hard cap on the number of nodes actually "
"expanded (default 200000). Remaining nodes are "
"emitted as stubs marked (node budget).")
parser.add_argument("--leaf-marker", type=str, default="(depth limit)",
help='marker appended to terminal nodes '
'(default: "(depth limit)"; use "" to suppress)')
parser.add_argument("--no-dedup", action="store_true",
help="expand every path in full instead of skipping "
"subtrees already shown elsewhere")
parser.add_argument("--no-classify", action="store_true",
help="do not annotate nodes with [def]/[thm]/[ax]")
parser.add_argument("--no-paths", action="store_true",
help="do not annotate nodes with their .lean file path")
group = parser.add_mutually_exclusive_group()
group.add_argument("--only-thms", action="store_true",
help="keep [thm] and [ax] nodes; prune [def] subtrees")
group.add_argument("--only-defs", action="store_true",
help="keep only [def] nodes; prune [thm]/[ax] subtrees")
parser.add_argument("--no-axioms", action="store_true",
help="drop [ax] nodes from the output "
"(compose with --only-thms for a pure theorem tree)")
args = parser.parse_args()

TARGET = args.target
MAX_DEPTH = None if args.depth <= 0 else args.depth
UNLIMITED = MAX_DEPTH is None
MAX_NODES = None if args.max_nodes <= 0 else args.max_nodes
OUT_FILE = args.out or (TARGET.replace(".", "_") + ".tree.txt")
CSV_FILE = args.csv
PATHS_FILE = args.paths
CLASSIFY = not args.no_classify
SHOW_PATHS = not args.no_paths
ONLY_THMS = args.only_thms
ONLY_DEFS = args.only_defs
NO_AXIOMS = args.no_axioms
LEAF_MARKER = args.leaf_marker
DEDUP = not args.no_dedup

if not os.path.isfile(CSV_FILE):
sys.exit(f"ERROR: CSV not found: {CSV_FILE}")

# ---------- column-name synonyms ----------
SRC_KEYS  = ("source", "src", "from", "decl", "declaration", "name",
"theorem", "node", "u", "parent")
TGT_KEYS  = ("target", "tgt", "to", "dest", "dependency", "dep",
"uses", "v", "child")
KIND_KEYS = ("kind", "class", "decl_kind", "declaration_kind",
"node_kind", "decltype", "decl_type", "is_def")
FILE_KEYS = ("file", "path", "file_path", "filename", "source_file",
"target_file", "src_file", "tgt_file", "mathlib_file",
"decl_file", "declaration_file", "location", "src_path",
"tgt_path", "source_path", "target_path")

# ---------- Lean kernel axioms ----------
# These four names are the primitive axioms of Lean 4's kernel.
#   propext          : (a <-> b) -> a = b
#   Quot.sound       : r a b -> Quot.mk r a = Quot.mk r b
#   Classical.choice : Nonempty a -> a
#   funext           : (forall x, f x = g x) -> f = g
# Mathlib conventionally lists all four when reporting #print axioms.
LEAN_AXIOMS = {
	"propext",
	"Quot.sound",
	"Classical.choice",
	"funext",
	# Alternate spellings that occasionally appear in CSVs
	"quot.sound",
	"classical.choice",
	"proof_irrel_heq",
	"Classical.propComplete",
}

# ---------- built-in heuristic namespace -> directory table ----------
NAMESPACE_PATHS = {
	"DifferentiableAt":             "Mathlib/Analysis/Calculus/Deriv/Basic.lean",
	"Differentiable":               "Mathlib/Analysis/Calculus/FDeriv/Basic.lean",
	"HasDerivAt":                   "Mathlib/Analysis/Calculus/Deriv/Basic.lean",
	"HasFDerivAt":                  "Mathlib/Analysis/Calculus/FDeriv/Basic.lean",
	"deriv":                        "Mathlib/Analysis/Calculus/Deriv/Basic.lean",
	"fderiv":                       "Mathlib/Analysis/Calculus/FDeriv/Basic.lean",
	"RCLike":                       "Mathlib/Analysis/RCLike/Basic.lean",
	"InnerProductSpace":            "Mathlib/Analysis/InnerProductSpace/Basic.lean",
	"ContinuousLinearMap":          "Mathlib/Topology/Algebra/Module/LinearMap.lean",
	"Continuous":                   "Mathlib/Topology/Defs/Basic.lean",
	"IsBoundedSMul":                "Mathlib/Analysis/NormedSpace/Basic.lean",
	"NormedAlgebra":                "Mathlib/Analysis/Normed/Algebra/Basic.lean",
	"intervalIntegral":             "Mathlib/MeasureTheory/Integral/IntervalIntegral/Basic.lean",
	"IntervalIntegrable":           "Mathlib/MeasureTheory/Integral/IntervalIntegral/Basic.lean",
	"MeasureTheory.MeasureSpace.volume":
	"Mathlib/MeasureTheory/MeasureSpace/Volume.lean",
	"MeasureTheory.MeasureSpace":   "Mathlib/MeasureTheory/MeasureSpace/Defs.lean",
	"MeasureTheory.Measure":        "Mathlib/MeasureTheory/Measure/Basic.lean",
	"MeasureTheory.IntegrableOn":   "Mathlib/MeasureTheory/Integral/OnFunctions.lean",
	"MeasureTheory.integral":       "Mathlib/MeasureTheory/Integral/Bochner/Basic.lean",
	"MeasureTheory":                "Mathlib/MeasureTheory",
	"Set.uIcc":                     "Mathlib/Data/Set/Intervals/UnorderedInterval.lean",
	"Set.Icc":                      "Mathlib/Data/Set/Intervals/Basic.lean",
	"Set.Ioo":                      "Mathlib/Data/Set/Intervals/Basic.lean",
	"Set.Ioc":                      "Mathlib/Data/Set/Intervals/Basic.lean",
	"Set.instMembership":           "Mathlib/Data/Set/Basic.lean",
	"Set.Mem":                      "Mathlib/Data/Set/Basic.lean",
	"Set":                          "Mathlib/Data/Set",
	"Real":                         "Mathlib/Data/Real/Basic.lean",
	"Nat":                          "Mathlib/Data/Nat",
	"Int":                          "Mathlib/Data/Int",
	"List":                         "Mathlib/Data/List",
	"Polynomial":                   "Mathlib/Algebra/Polynomial",
	"PseudoMetricSpace":            "Mathlib/Topology/MetricSpace/PseudoMetricSpace.lean",
	"UniformSpace":                 "Mathlib/Topology/UniformSpace/Basic.lean",
	"Bornology":                    "Mathlib/Topology/Bornology/Basic.lean",
	"CompleteSpace":                "Mathlib/Topology/Algebra/UniformGroup.lean",
	"TopologicalSpace":             "Mathlib/Topology/Defs/Basic.lean",
	"NormedAddCommGroup":           "Mathlib/Analysis/Normed/Group/Basic.lean",
	"NormedAddGroup":               "Mathlib/Analysis/Normed/Group/Basic.lean",
	"NormedCommRing":               "Mathlib/Analysis/Normed/Ring/Basic.lean",
	"NormedRing":                   "Mathlib/Analysis/Normed/Ring/Basic.lean",
	"NonUnitalNormedCommRing":      "Mathlib/Analysis/Normed/Ring/Basic.lean",
	"NonUnitalNormedRing":          "Mathlib/Analysis/Normed/Ring/Basic.lean",
	"SeminormedAddCommGroup":       "Mathlib/Analysis/Normed/Group/Basic.lean",
	"SeminormedRing":               "Mathlib/Analysis/Normed/Ring/Basic.lean",
	"SeminormedCommRing":           "Mathlib/Analysis/Normed/Ring/Basic.lean",
	"NonUnitalSeminormedCommRing":  "Mathlib/Analysis/Normed/Ring/Basic.lean",
	"NonUnitalSeminormedRing":      "Mathlib/Analysis/Normed/Ring/Basic.lean",
	"DenselyNormedField":           "Mathlib/Analysis/Normed/Field/Basic.lean",
	"NontriviallyNormedField":      "Mathlib/Analysis/Normed/Field/Basic.lean",
	"NormedField":                  "Mathlib/Analysis/Normed/Field/Basic.lean",
	"NormedSpace":                  "Mathlib/Analysis/NormedSpace/Basic.lean",
	"dist":                         "Mathlib/Topology/MetricSpace/PseudoMetricSpace.lean",
	"Dist":                         "Mathlib/Topology/MetricSpace/PseudoMetricSpace.lean",
	"enorm":                        "Mathlib/Analysis/Normed/Group/Basic.lean",
	"AddGroup":                     "Mathlib/Algebra/Group/Defs.lean",
	"AddCommGroup":                 "Mathlib/Algebra/Group/Defs.lean",
	"AddMonoid":                    "Mathlib/Algebra/Group/Defs.lean",
	"AddZero":                      "Mathlib/Algebra/Group/Defs.lean",
	"AddZeroClass":                 "Mathlib/Algebra/Group/Defs.lean",
	"SubNegMonoid":                 "Mathlib/Algebra/Group/Defs.lean",
	"SubNegZeroMonoid":             "Mathlib/Algebra/Group/Defs.lean",
	"SubtractionCommMonoid":        "Mathlib/Algebra/Group/Defs.lean",
	"MonoidWithZero":               "Mathlib/Algebra/GroupWithZero/Defs.lean",
	"MulZeroClass":                 "Mathlib/Algebra/GroupWithZero/Defs.lean",
	"NegZeroClass":                 "Mathlib/Algebra/Group/Defs.lean",
	"Ring":                         "Mathlib/Algebra/Ring/Defs.lean",
	"CommRing":                     "Mathlib/Algebra/Ring/Defs.lean",
	"Semiring":                     "Mathlib/Algebra/Ring/Defs.lean",
	"Semifield":                    "Mathlib/Algebra/Field/Defs.lean",
	"DivisionSemiring":             "Mathlib/Algebra/Field/Defs.lean",
	"Field":                        "Mathlib/Algebra/Field/Defs.lean",
	"Module":                       "Mathlib/Algebra/Module/Defs.lean",
	"DistribMulAction":             "Mathlib/Algebra/Module/Defs.lean",
	"DistribSMul":                  "Mathlib/Algebra/Module/Defs.lean",
	"SMulZeroClass":                "Mathlib/Algebra/Module/Defs.lean",
	"Lattice":                      "Mathlib/Order/Lattice.lean",
	"DistribLattice":               "Mathlib/Order/Lattice.lean",
	"SemilatticeInf":               "Mathlib/Order/Lattice.lean",
	"SemilatticeSup":               "Mathlib/Order/Lattice.lean",
	"PartialOrder":                 "Mathlib/Order/Defs.lean",
	"Preorder":                     "Mathlib/Order/Defs.lean",
	"LinearOrder":                  "Mathlib/Order/Defs.lean",
	"HSub":                         "Mathlib/Algebra/Notation.lean",
	"instHSub":                     "Mathlib/Algebra/Notation.lean",
	"Subtype":                      "Mathlib/Data/Subtype/Basic.lean",
	"Zero":                         "Mathlib/Algebra/Group/Defs.lean",
	"OfNat":                        "Mathlib/Algebra/Group/Defs.lean",
}

LEAN_CORE = {
	"forall_congr", "implies_congr", "inferInstance", "outParam",
	"OfNat", "LE", "LT", "Max", "Min", "Membership", "Zero", "One",
	"HAdd", "HMul", "HPow", "HSub", "HDiv", "HMod", "HAnd", "HOr",
	"HXor", "HShiftLeft", "HShiftRight", "HAppend", "HEq", "HImp",
	"HIff", "HNot", "HSMul", "Decidable", "Nonempty", "PLift", "ULift",
	"Quot", "Quotient", "Classical", "Eq",
	"Subsingleton", "Inhabited", "Empty", "PUnit",
	"Unit", "Bool", "Option", "Sum", "Prod", "Sigma", "PSigma",
	"eq_self", "of_eq_true", "le_total",
	"min_comm", "min_eq_left", "min_eq_right",
	"sup_eq_left", "sup_eq_right",
}

# ---------- CSV helpers ----------
def _index_of(header_lower, keys, default=None):
for i, c in enumerate(header_lower):
if c in keys:
return i
return default

def _indices_of(header_lower, keys):
return [i for i, c in enumerate(header_lower) if c in keys]

def load_paths_file(path):
m = {}
if not path or not os.path.isfile(path):
return m
with open(path, "r", encoding="utf-8") as f:
first = f.readline().rstrip("\n")
cols = [c.strip().lower() for c in first.split(",")]
if "name" in cols and ("file" in cols or "path" in cols):
name_idx = cols.index("name")
file_idx = cols.index("file") if "file" in cols else cols.index("path")
else:
name_idx, file_idx = 0, 1
parts = first.split(",")
if len(parts) > 1:
m[parts[0].strip()] = parts[1].strip()
for line in f:
parts = line.rstrip("\n").split(",")
if len(parts) <= max(name_idx, file_idx):
continue
m[parts[name_idx].strip()] = parts[file_idx].strip()
return m

# ---------- load edges ----------
print(f"Loading edges from {CSV_FILE} ...")
adj = defaultdict(list)
known_nodes = set()
kind_of = {}
csv_paths = {}

with open(CSV_FILE, "r", encoding="utf-8") as f:
header_line = f.readline().rstrip("\n")
header_cols = [c.strip() for c in header_line.split(",")]
header_lower = [c.lower() for c in header_cols]

src_idx  = _index_of(header_lower, SRC_KEYS, 0)
tgt_idx  = _index_of(header_lower, TGT_KEYS, 1)
kind_idx = _index_of(header_lower, KIND_KEYS, None)
file_idxs = _indices_of(header_lower, FILE_KEYS)

print(f"  Header: {header_cols}")
print(f"  -> source column: '{header_cols[src_idx]}' (index {src_idx})")
print(f"  -> target column: '{header_cols[tgt_idx]}' (index {tgt_idx})")
if kind_idx is not None:
print(f"  -> kind column:   '{header_cols[kind_idx]}' (index {kind_idx})")
else:
print("  -> kind column:   not present (heuristic will be used)")
if file_idxs:
names = ", ".join(f"'{header_cols[i]}'" for i in file_idxs)
print(f"  -> file columns:  {names}")
else:
print("  -> file columns:  none detected")

for line in f:
parts = line.rstrip("\n").split(",")
if len(parts) <= max(src_idx, tgt_idx):
continue
src = parts[src_idx].strip()
tgt = parts[tgt_idx].strip()
adj[src].append(tgt)
known_nodes.add(src)
known_nodes.add(tgt)

if kind_idx is not None and len(parts) > kind_idx:
raw = parts[kind_idx].strip().lower()
if raw in ("def", "definition", "abbrev", "instance", "structure"):
kind_of[src] = "def"
elif raw in ("thm", "theorem", "lemma", "prop", "proposition",
"corollary"):
kind_of[src] = "thm"
elif raw in ("ax", "axiom"):
kind_of[src] = "ax"

for fi in file_idxs:
if len(parts) <= fi:
continue
val = parts[fi].strip()
if not val:
continue
col_name = header_lower[fi]
if "target" in col_name:
csv_paths.setdefault(tgt, val)
else:
csv_paths.setdefault(src, val)

edges_total = sum(len(v) for v in adj.values())
print(f"  {len(adj):,} source nodes, {edges_total:,} edges")
print(f"  {len(kind_of):,} nodes carry kind info from CSV")
print(f"  {len(csv_paths):,} nodes carry exact file paths from CSV")

PATHS_MAP = load_paths_file(PATHS_FILE)
for k, v in PATHS_MAP.items():
csv_paths[k] = v
if PATHS_MAP:
print(f"  {len(PATHS_MAP):,} additional mappings loaded from {PATHS_FILE}")

if TARGET not in known_nodes:
print(f"WARNING: '{TARGET}' not found in graph - tree may be empty.")

# ---------- path resolution ----------
def path_of(name):
if name in csv_paths:
return csv_paths[name]
best = None
for k in csv_paths:
if name.startswith(k + ".") and (best is None or len(k) > len(best)):
best = k
if best is not None:
return csv_paths[best]
# axioms live in Lean core / Mathlib's Logic.lean
if name in LEAN_AXIOMS:
if name in ("Classical.choice", "classical.choice",
"Classical.propComplete"):
return "Mathlib/Logic/Choice.lean (Lean core)"
return "Mathlib/Logic/Basic.lean (Lean core)"
best = None
for k in NAMESPACE_PATHS:
if (name == k or name.startswith(k + ".")) and \
(best is None or len(k) > len(best)):
best = k
if best is not None:
return NAMESPACE_PATHS[best]
if "." in name:
top = name.split(".", 1)[0]
if top in NAMESPACE_PATHS:
return NAMESPACE_PATHS[top]
if "." not in name or name.split(".", 1)[0] in LEAN_CORE:
return "Lean core (not in Mathlib)"
return f"Mathlib/**/{name.split('.', 1)[0]}/**"

# ---------- classification ----------
_INHERITANCE_RE = re.compile(r"\.to[A-Z]")

def classify_heuristic(name):
# axiom check comes first
if name in LEAN_AXIOMS:
return "ax"
if not name:
return "def"
last = name.rsplit(".", 1)[-1]
if last.startswith("inst") or ".inst" in name:
return "def"
if _INHERITANCE_RE.search(name):
return "def"
if last and last[0].isupper():
return "def"
if "_" in last:
return "thm"
return "def"

def node_kind(name):
if name in kind_of:
return kind_of[name]
return classify_heuristic(name)

def tag(name):
"""Class tag only; the path is emitted on its own line."""
if not CLASSIFY:
return ""
return f"  [{node_kind(name)}]"

def keep(name):
"""Return True if this node should be kept under the active filters."""
k = node_kind(name)
if NO_AXIOMS and k == "ax":
return False
if not (ONLY_THMS or ONLY_DEFS):
return True
if ONLY_THMS:
# keep theorems AND axioms: axioms are the foundation of proofs
return k in ("thm", "ax")
if ONLY_DEFS:
return k == "def"
return True

# ---------- tree data structure ----------
class TNode:
__slots__ = ("name", "children", "marker", "is_cycle")
def __init__(self, name):
self.name = name
self.children = []
self.marker = ""
self.is_cycle = False

_nodes_created = 0  # respect --max-nodes

def build_tree(node, depth=0, visited=None, seen_on_path=None):
global _nodes_created
if visited is None:
visited = set()
if seen_on_path is None:
seen_on_path = set()

n = TNode(node)

if node in seen_on_path:
n.is_cycle = True
return n
if DEDUP and node in visited:
n.marker = "(*)"
return n

if MAX_NODES is not None and _nodes_created >= MAX_NODES:
n.marker = "(node budget)"
return n
_nodes_created += 1

if DEDUP:
visited.add(node)
seen_on_path.add(node)

if MAX_DEPTH is not None and depth >= MAX_DEPTH:
n.marker = "(depth limit)"
seen_on_path.discard(node)
return n

deps = adj.get(node, [])
if not deps:
n.marker = LEAF_MARKER if LEAF_MARKER else ""
seen_on_path.discard(node)
return n

for dep in sorted(set(deps)):
n.children.append(build_tree(dep, depth + 1, visited, seen_on_path))

seen_on_path.discard(node)
return n

def prune_tree(node):
if not keep(node.name):
return None
new = TNode(node.name)
new.marker = node.marker
new.is_cycle = node.is_cycle
for child in node.children:
pc = prune_tree(child)
if pc is not None:
new.children.append(pc)
return new

def render_struct(node, prefix="", is_root=True, is_last=True):
"""Two lines per node: declaration + tag, then `@ PATH`.

The second line reuses the same ancestor `prefix` as the first
so the `@` sits flush under the branch marker.
"""
lines = []
display = f"(cycle: {node.name})" if node.is_cycle else node.name

# ---- name line ------------------------------------------------
if is_root:
name_line = display + tag(node.name)
else:
name_line = (prefix
+ ("`-- " if is_last else "|-- ")
+ display
+ tag(node.name))
if node.marker:
name_line += f"   {node.marker}"
lines.append(name_line)

# ---- path line ------------------------------------------------
if SHOW_PATHS:
lines.append(prefix + "@ " + path_of(node.name))

# ---- children -------------------------------------------------
child_prefix = "" if is_root else prefix + ("    " if is_last else "|   ")
n = len(node.children)
for i, child in enumerate(node.children):
lines.extend(render_struct(
child, prefix=child_prefix, is_root=False,
is_last=(i == n - 1)))
return lines

def count_nodes(node):
total = 1
for c in node.children:
total += count_nodes(c)
return total

# ---------- main ----------
print(f"Generating dependency tree for: {TARGET}")
if UNLIMITED:
print("  --depth 0 (unlimited expansion)")
if MAX_NODES is not None:
print(f"  --max-nodes {MAX_NODES:,}")
if not DEDUP:
print("  --no-dedup (subtrees re-expanded at every occurrence)")

root = build_tree(TARGET, depth=0)
raw_size = count_nodes(root)
budget_hit = MAX_NODES is not None and _nodes_created >= MAX_NODES
print(f"  raw tree nodes: {raw_size:,}")
if budget_hit:
print(f"  WARNING: node budget ({MAX_NODES:,}) reached; tree was truncated.")

if not DEDUP and raw_size > 200_000:
print(f"  WARNING: --no-dedup produced {raw_size:,} nodes; "
f"the output file may be very large.")

if ONLY_THMS or ONLY_DEFS or NO_AXIOMS:
pruned = prune_tree(root)
if pruned is None:
pruned = TNode(root.name)
pruned.marker = "(filtered out: target does not match --only-* filter)"
root = pruned

lines = render_struct(root)

n_def = sum(1 for ln in lines if "[def]" in ln)
n_thm = sum(1 for ln in lines if "[thm]" in ln)
n_ax  = sum(1 for ln in lines if "[ax]"  in ln)
n_leaf = sum(1 for ln in lines if LEAF_MARKER and LEAF_MARKER in ln)

header = [
"=" * 60,
f"Dependency tree for: {TARGET}",
f"Source: {CSV_FILE}",
f"Max depth: {'unlimited' if UNLIMITED else MAX_DEPTH}",
"(*) = already expanded elsewhere in this tree",
]
if CLASSIFY:
header.append("[def] = definition / typeclass instance")
header.append("[thm] = theorem / lemma")
header.append("[ax]  = Lean kernel axiom "
"(propext, funext, Quot.sound, Classical.choice)")
if SHOW_PATHS:
if csv_paths:
header.append("@ PATH = on the line below each declaration, "
"exact from mathlib_edges.csv (or --paths) "
"where available; heuristic otherwise")
else:
header.append("@ PATH = on the line below each declaration, "
"heuristic from namespace only "
"(no path columns in CSV, no --paths file)")
if LEAF_MARKER:
header.append(f"{LEAF_MARKER} = node has no further dependencies "
f"in the CSV")
if not DEDUP:
header.append("Dedup: disabled (subtrees re-expanded)")
if ONLY_THMS:
header.append("Filter: --only-thms "
"(pruned [def] subtrees; [ax] retained)")
elif ONLY_DEFS:
header.append("Filter: --only-defs (pruned [thm]/[ax] subtrees)")
if NO_AXIOMS:
header.append("Filter: --no-axioms (dropped [ax] nodes)")
if MAX_NODES is not None:
header.append(f"Node budget: {MAX_NODES:,}"
+ ("  [HIT - output truncated]"
if budget_hit else ""))
header.append("=" * 60)
header.append("")

with open(OUT_FILE, "w", encoding="utf-8") as f:
f.write("\n".join(header + lines))
f.write("\n")

print(f"Done. Written to: {OUT_FILE}")
print(f"  {len(lines):,} lines")
if CLASSIFY:
print(f"  {n_def:,} def, {n_thm:,} thm, {n_ax:,} ax")
if LEAF_MARKER:
print(f"  {n_leaf:,} terminal nodes ({LEAF_MARKER})")
\end{verbatimbox}
\setstretch{1}

\section{Tables}

\label{sec:tables}

\subsection{Timeline}

\label{sec:timeline}

\noindent Key dates in the formalisation of FTC \cite{mathlib-stats,
mathlib4}.\text{ \ }\newline
\text{ \ }\newline
\setstretch{1.3} $%
\begin{array}{|T|S|L|K|}
\hline
\hypertarget{tl:2013}{\textbf{2013}} &
first Lean~2 formal\allowbreak isation &
\textsf{integral\allowbreak\_deriv\allowbreak\_eq\allowbreak\_sub} \textsf{(Lean 2 name)} &
\textsf{Hoel et al.} \\ \hline
\hypertarget{tl:2017}{\textbf{2017}} &
Lean~3 port &
\textsf{integral\allowbreak\_eq\allowbreak\_sub\allowbreak\_of\allowbreak\_has\allowbreak Deriv\allowbreak At} &
\textsf{math\allowbreak lib3} \\ 
\end{array}$

\noindent$%
\begin{array}{|T|S|L|K|}
\hypertarget{tl:2021}{\textbf{2021}} &
inter\allowbreak val integ\allowbreak ral rewrit\allowbreak ten &
\textsf{FTC\allowbreak Filter} \textsf{introduced} &
\textsf{math\allowbreak lib3 PR} \\ \hline
\hypertarget{tl:2022}{\textbf{2022}} &
Math\allowbreak lib~4 port &
mathlib4 naming conven\allowbreak tion &
\textsf{math\allowbreak lib4} \\ 
\end{array}$

\noindent$%
\begin{array}{|T|S|L|K|}
\hline
\hypertarget{tl:2024}{\textbf{2024}} &
gener\allowbreak al Lit\allowbreak tle-$o$ FTC-2 stabil\allowbreak ised &
\textsf{measure\allowbreak\_integral\allowbreak\_sub\allowbreak\_integral\allowbreak\_sub\allowbreak\_linear\allowbreak\_is\allowbreak LittleO\allowbreak\_of\allowbreak\_tendsto\allowbreak\_ae} &
\textsf{math\allowbreak lib4} \\ \hline
\end{array}$ \setstretch{1}

\subsection{Classifications}

\label{sec:classifications}

\noindent Every node of the depth-5 tree falls into one of four classes.
\text{ \  \ }\newline
\text{ \  \ }\newline
$%
\begin{array}{|T|S|L|K|}
\hline
\hypertarget{cl:ax}{\textbf{Ax}} &
2 visible at depth~5 in excerpt; 3 reached overall &
\textsf{funext},\;\textsf{Quot\allowbreak.sound};\;
\textsf{propext} (supplement) &
\textsf{ker\allowbreak nel ax\allowbreak ioms} \\ \hline
\hypertarget{cl:def}{\textbf{Def}} &
domin\allowbreak ant class &
\textsf{Interval\allowbreak Integrable},\;\textsf{FTC\allowbreak Filter},\;\textsf{Normed\allowbreak AddCommGroup} &
\textsf{type\allowbreak class clo\allowbreak sure} \\ \hline
\end{array}$ \text{ \  \ }\newline\noindent $%
\begin{array}{|T|S|L|K|}
\hypertarget{cl:lem}{\textbf{Lem}} &
many tech\allowbreak nical steps &
\textsf{Is\allowbreak LittleO.add\allowbreak\_add},\;\textsf{Filter.mp\allowbreak\_mem} &
\textsf{tech\allowbreak nical steps}
\\ \hline
\hypertarget{cl:thm}{\textbf{Thm}} &
$\approx 25$ substantive nodes & 
\textsf{measure\allowbreak\_integral\allowbreak\_sub\allowbreak\_linear\allowbreak\_is\allowbreak LittleO\allowbreak\_of\allowbreak\_tendsto\allowbreak\_ae} &
\textsf{mathe\allowbreak matical con\allowbreak tent} 
 \\ \hline
\end{array}$

\noindent \newline

\subsection{Applications}

\label{sec:applications} \setstretch{1.5} \noindent The general Little-$o$
form of FTC-2 above underpins: \setstretch{1.5} \text{ \ }\newline
\text{ \ }

\noindent$%
\begin{array}{|T|S|L|K|}
\hline
\hypertarget{ap:cov}{\textbf{App\,1}} &
change of vari\allowbreak ables &
\textsf{integral\allowbreak\_comp\allowbreak\_mul\allowbreak\_deriv} &
\textsf{Change of Vari\allowbreak ables} \\ \hline
\hypertarget{ap:barrow}{\textbf{App\,2}} &
inverse integr\allowbreak ation / different\allowbreak iation &
\textsf{integral\allowbreak\_eq\allowbreak\_sub\allowbreak\_of\allowbreak\_has\allowbreak Deriv\allowbreak At} &
\textsf{Barrow's Rule} \\
\end{array}$

\noindent$%
\begin{array}{|T|S|L|K|}
\hline
\hypertarget{ap:Wiedijk}{\textbf{App\,3}} &
FTC as named target &
\textsf{integral\allowbreak\_deriv\allowbreak\_eq\allowbreak\_sub} &
\textsf{Wied\allowbreak ijk Thm 14} \\ \hline
\hypertarget{ap:circle}{\textbf{App\,4}} &
computes $\pi$ &
\textsf{integral\allowbreak\_sqrt\allowbreak\_one\allowbreak\_sub\allowbreak\_sq} &
\textsf{Area of the Cir\allowbreak cle} \\ 
\end{array}$

\noindent$%
\begin{array}{|T|S|L|K|}
\hline
\hypertarget{ap:bochner}{\textbf{App\,5}} &
$E$-valued Boch\allowbreak ner general\allowbreak isation &
\textsf{MeasureTheory.integral}\;$[Complete\allowbreak Space\;E]$ &
\textsf{Bochn\allowbreak er Integ\allowbreak ral FTC} \\ \hline
\end{array}$ \newpage

\section{Conclusion}

\label{sec:conclusion} \noindent \setstretch{1.3} The Fundamental Theorem of
Calculus in Mathlib is not a single theorem but a small hierarchy whose most
general member is the Little-$o$ statement \textsf{interval\allowbreak
Integral.measure\allowbreak
\_integral\allowbreak\_sub\allowbreak\_integral\allowbreak\_sub\allowbreak
\_linear\allowbreak\_is\allowbreak LittleO\allowbreak\_of\allowbreak
\_tendsto\allowbreak\_ae}. From it one derives the classical derivative form 
\textsf{interval\allowbreak Integral.integral\allowbreak\_eq\allowbreak
\_sub\allowbreak\_of\allowbreak\_has\allowbreak Deriv\allowbreak At} and the
one-sided variants \textsf{\_left} and \textsf{\_right}. The dependency tree
in Section~\ref{sec:tree} makes this structure explicit.

\noindent The depth-5 closure of the one-sided theorem contains on the order
of $10^{5}$ nodes. Two of the three kernel axioms reachable in the closure
--- \textsf{funext} and \textsf{Quot\allowbreak.sound} --- are visible in
the printed excerpt; \textsf{propext} appears immediately beyond it on the 
\textsf{Asymptotics.IsLittleO.sub} branch. \textsf{Classical
\allowbreak.choice} is not reached. The same methodology --- extract the
dependency tree, cross-check signatures against the documentation \cite%
{mathlib4-docs}, and inspect the kernel-axiom closure --- applies to any
other theorem in the library \cite{mathlib4, mathlibgraph2026}. The Lean
skeleton of Section~\ref{sec:lean-skeleton} has been type-checked against
Mathlib commit \texttt{0c1d5a7e}{} and introduces no new axiom dependencies
beyond those already present in the one-sided statement.

\begin{thebibliography}{99}
\bibitem{close2026axiom} Larsen James Close. \newblock
The Axiom Profile of Mathlib. \newblock
Zenodo, 2026. \newblock
DOI: \texttt{10.5281/zenodo.20710920}. \newblock
\url{https://zenodo.org/records/20710921}.

\medskip

\bibitem{deepseek-chat} DeepSeek. \newblock
DeepSeek Chat. \newblock
Conversational assistant used to assemble the exposition, cross-check the 
statement signatures against the Mathlib source, and prepare the LaTeX 
source of this paper. \newblock
\url{https://chat.deepseek.com}.

\medskip

\bibitem{mathlib-stats} Lean community. \newblock
Mathlib statistics. \newblock
\url{https://leanprover-community.github.io/mathlib_stats.html}.

\medskip

\bibitem{mathlib4} The mathlib Community. \newblock
The Lean mathematical library. \newblock
Source repository, GitHub. \newblock
\url{https://github.com/leanprover-community/mathlib4}.

\medskip

\bibitem{mathlib4-docs} The mathlib Community. \newblock
Mathlib4 documentation browser. \newblock
Online reference for every declaration in Mathlib; used to locate the 
signatures and dependency edges of the theorems in this paper, and 
recommended as the entry point for analogous dependency-guided expositions 
of other theorems in the same library. \newblock
\url{https://leanprover-community.github.io/mathlib4_docs/}.

\medskip

\bibitem{mathlib4-fundthm} The mathlib Community. \newblock
\texttt{Mathlib/MeasureTheory/Integral/IntervalIntegral/FundThmCalculus.lean}%
. \newblock
Mathlib4 source file containing the two-sided Little-$o$ FTC-2 and the 
one-sided reduction used in Section~\ref{sec:proofs} and  Section~\ref%
{sec:scripts}. \newblock
Pinned to commit \texttt{0c1d5a7e}{}. \newblock
\url{https://github.com/leanprover-community/mathlib4/blob/%
\mathlibsharaw/Mathlib/MeasureTheory/Integral/IntervalIntegral/FundThmCalculus.lean}%
.

\medskip

\bibitem{mathlibgraph} Xinze Li and Patrick Shafto. \newblock
The Multinetwork of Mathlib. \newblock
NeurIPS 2026 Datasets and Benchmarks Track, 2026. Under review. \newblock
Dataset file: \texttt{mathlib\_edges.csv}. \newblock
\url{https://huggingface.co/datasets/MathNetwork/MathlibGraph/resolve/main/mathlib_edges.csv}%
.

\medskip

\bibitem{mathlibgraph2026} Xinze Li and Patrick Shafto. \newblock
MathlibGraph: The Multinetwork of Mathlib. \newblock
Hugging Face dataset, 2026. \newblock
\url{https://huggingface.co/datasets/MathNetwork/MathlibGraph}.

\medskip

\bibitem{network-structure-mathlib} Xinze Li, Nanyun Peng, Simone Severini, 
and Patrick Shafto. \newblock
The Network Structure of Mathlib. \newblock
arXiv:2604.24797, 2026. \newblock
\url{https://arxiv.org/abs/2604.24797}.
\end{thebibliography}

\end{document}
